% Chapter 1: Graham Positivity of Triple Schubert Calculus
% 基于 Yibo Gao, Rui Xiong (2025)
% arXiv: 2506.09421

\chapter{Graham Positivity of Triple Schubert Calculus}
\label{chap:TriplePositivity}

\section{预备知识}
\label{sec:Preliminaries}

在本节中，我们介绍本文所需的基本定义.

\subsection{旗流形}

\begin{Definition}[旗流形（Flag variety）]
\label{def:FlagVariety}
设 $V$ 为 $n$ 维复向量空间. 旗流形 $Fl_n(\mathbb{C})$ 是 $V$ 中所有 flags 的集合，其中一个 flag 是一系列子空间
\begin{equation}
0 = V_0 \subset V_1 \subset V_2 \subset \cdots \subset V_n = V
\end{equation}
其中 $\dim V_i = i$. 旗流形是一个光滑的复代数簇，维度为 $\binom{n}{2}$.
\end{Definition}

旗流形是 Schubert 理论的舞台. Schubert 胞腔是旗流形的重要子簇.

\subsection{对称群与排列}

\begin{Definition}[对称群（Symmetric group）]
\label{def:SymmetricGroup}
\textbf{\underline{$n$ 阶对称群}} $S_n$ 是 $\{1, 2, \ldots, n\}$ 的所有置换构成的群. 一个置换 $w$ 可以记为一维数组 $w = (w(1), w(2), \ldots, w(n))$.

排列 $w$ 的\textbf{\underline{逆序数（inversion number）}} $\ell(w)$ 定义为满足 $1 \leq i < j \leq n$ 且 $w(i) > w(j)$ 的对 $(i, j)$ 的个数，也称为 $w$ 的长度.

在 $S_n$ 中，长度最大的置换称为\textbf{\underline{最长置换}} $w_0$，它将 $\{1,2,\ldots,n\}$ 逆序排列：
\[
w_0 = (n, n-1, \ldots, 2, 1)
\]
其长度为 $\ell(w_0) = \binom{n}{2}$.

$S_\infty$ 表示所有有限置换的并集，即 $S_\infty = \bigcup_{n=1}^\infty S_n$\footnote{问：$S_\infty$ 定义中的 $\bigcup$（并集记号）是什么意思？见附录 \cref{sec:SInfinity}。}，可以理解为可数无限集 $\{1, 2, \ldots\}$ 上的对称群。例如，$w = (3,1,4,2,5,\ldots) \in S_\infty$ 表示一个有限支撑的置换（只有前4位发生交换）。

\textbf{排列记号的省略约定}：当我们写 $w = 21$ 时，表示的是将 $1$ 和 $2$ 交换、$3$ 保持不变的置换，即一行数组 $(2,1,3)$。换言之，\textbf{若排列的后 $n-k$ 个位置不变，则这些位置可以省略}。例如：
\begin{itemize}
\item 在 $S_3$ 中：$21$ 表示 $(2,1,3)$，即 $w(1)=2, w(2)=1, w(3)=3$
\item 在 $S_n$ 中：$21$ 表示 $(2,1,3,4,\ldots,n)$
\item $132$ 表示 $(1,3,2)$，即交换 $2$ 和 $3$，$1$ 不变
\end{itemize}
这个约定在上下文明确时简化记号，但需注意避免歧义。\footnote{当需明确时，写完整的一行数组如 $(2,1,3)$ 或用置换符号如 $s_1$。}
\end{Definition}

\begin{Definition}[Bruhat 序（Bruhat order）]
\label{def:BruhatOrder}
设 $u, v \in S_n$，称 $u \leq v$（在 Bruhat 序下），如果 $u$ 可以通过一系列基本置换的乘积得到 $v$，且这些基本置换对应的指标序列是 $v$ 的某个约化表示的子序列.

等价的组合描述：$u \leq v$ 当且仅当对所有 $1 \leq k \leq n$ 和 $j$，有
\[
\#\{i \leq k : u(i) \leq j\} \leq \#\{i \leq k : v(i) \leq j\}.
\]
\end{Definition}

\subsection{代数几何中的基本相交概念}
\label{sec:IntersectionConcepts}

在我们讨论旗流形上 Schubert 胞腔的相交时，两个术语至关重要——它们出现在 Lemma 2.5（横截相交引理）中：

\begin{Definition}[Proper（正常/固有）交集]
\label{def:ProperIntersection}
设 $X, Y \subset Z$ 是复代数簇的闭子簇。称交集 $X \cap Y$ 是 \textbf{proper}（正常的/固有的），如果交集 $X \cap Y$ 是紧的（compact），或者在相对版本中，映射 $X \cap Y \to Z$ 是 proper 态射。

直观理解：proper 意味着交集``行为良好''——不会出现无穷多点，也不会在极限处丢失点。
\end{Definition}

\begin{Definition}[Transverse（横截）交集]
\label{def:TransverseIntersection}
设 $X, Y \subset Z$ 是光滑代数簇的闭子簇。称 $X$ 与 $Y$ 在点 $p \in X \cap Y$ 处\textbf{横截（transverse）}相交，如果
\[
T_p X + T_p Y = T_p Z,
\]
即 $X$ 和 $Y$ 在 $p$ 处的切空间张成 $Z$ 在 $p$ 处的整个切空间。

等价描述：横截相交意味着两子簇在交点处``横切''——没有相切、重合或包含关系。

\textbf{横截相交的性质}：
\begin{itemize}
\item 横截交集点是孤立的（discrete）——交点数局部是常数；
\item 横截相交在形变下保持稳定——稍微移动子簇不会改变交点个数；
\item 横截相交的点个数是良定义的整数（良定性）。
\end{itemize}

\textbf{非横截相交的例子}：两条平面曲线在原点相切——它们有无穷多个交点（整个切点邻域），而不是孤立的点。
\end{Definition}

\textbf{为什么 Lemma 2.5 关注这两个性质？} Proper 保证交集是紧的（有限多个点），transverse 保证每个交点都是良定的（横截点）。两者结合使得在等变量子上同调中可以定义明确的相交类，从而得到非负的展开系数。

\begin{Example}[Bruhat 序（Bruhat order）的例子]
在 $S_3$ 中，$123$ 是最小元，$321$ 是最大元。S$_3$ 的 Bruhat order 是偏序（partial order），其 Hasse 图为：
\[
\begin{tikzpicture}[node distance=1.5cm]
\node (321) at (0,4) {$321$};
\node (312) at (2,2.5) {$312$};
\node (231) at (-2,2.5) {$231$};
\node (132) at (-1,0.5) {$132$};
\node (213) at (1,0.5) {$213$};
\node (123) at (0,-0.5) {$123$};

\draw (123) -- (132);
\draw (123) -- (213);
\draw (132) -- (231);
\draw (213) -- (312);
\draw (231) -- (321);
\draw (312) -- (321);
\end{tikzpicture}
\]
即
\[
123 < 213, \quad 123 < 132, \quad 213 < 321, \quad 132 < 321,
\]
且 $213$ 与 $132$，以及 $231$ 与 $312$ 在 Bruhat 序下\textbf{不可比较（incomparable）}。
\end{Example}

\begin{Definition}[Descents（下降点）]
\label{def:Descents}
置换 $w \in S_n$ 的 \textbf{\underline{descents 集合}}定义为：
$$\mathrm{Des}(w) = \{i \in \{1, \ldots, n-1\} : w(i) > w(i+1)\}$$
即满足 $w(i) > w(i+1)$ 的位置 $i$ 的集合. 每个 descent 将置换分为若干块.

例如：$w = 35142$（即 $(3,5,1,4,2)$），有 $w(1)=3 > w(2)=5$ 不成立，$w(2)=5 > w(3)=1$ 成立，$w(3)=1 > w(4)=4$ 不成立，$w(4)=4 > w(5)=2$ 成立，故 $\mathrm{Des}(w) = \{2, 4\}$.
\end{Definition}

\begin{Definition}[Separated Descents]
\label{def:SeparatedDescents}
设 $u, v \in S_\infty$。若 $\max \operatorname{Des}(u) \leq \min \operatorname{Des}(v)$，则 $(u, v)$ 满足 separated descents 条件。

即 $u$ 的最大下降点不超过 $v$ 的最小下降点，保证了两者的下降区域不重叠。
\end{Definition}

\subsection{除差算子}
\label{sec:DividedDifference}

\begin{Definition}[除差算子（Divided difference operator）]
\label{def:DividedDifference}
对于 $i = 1, \ldots, n-1$，除差算子 $\partial_i$ 定义为
\begin{equation}
\partial_i(f) = \frac{f(\ldots, x_i, x_{i+1}, \ldots) - f(\ldots, x_{i+1}, x_i, \ldots)}{x_i - x_{i+1}}.
\end{equation}
更一般地，对于置换 $w$，定义其对应的除差算子 $\partial_w$ 为适当组合的 $\partial_i$.
\end{Definition}

\begin{Example}[除差算子的例子]
\label{ex:DividedDifference}
在 $S_3$ 中（变量为 $x_1, x_2, x_3$）：
\begin{itemize}
\item $\partial_1(x_1) = \frac{x_1 - x_2}{x_1 - x_2} = 1$。
\item $\partial_1(x_2) = \frac{x_2 - x_1}{x_1 - x_2} = -1$。
\item $\partial_1(x_1^2) = \frac{x_1^2 - x_2^2}{x_1 - x_2} = x_1 + x_2$。
\item $\partial_1(x_1 x_2) = \frac{x_1 x_2 - x_2 x_1}{x_1 - x_2} = 0$。
\item $\partial_1(x_2^2) = \frac{x_2^2 - x_1^2}{x_1 - x_2} = -(x_1 + x_2)$。
\item $\partial_1 \partial_2(x_1) = \partial_1(1) = 0$（因为 $\partial_2$ 作用于 $x_1$ 得 1，再 $\partial_1$ 作用于常数得 0）。
\item $\partial_2(x_2) = 1$（$\partial_2$ 交换 $x_2$ 和 $x_3$）。
\end{itemize}
除差算子的作用是取对称化：$\partial_i$ 将多项式变为关于 $x_i, x_{i+1}$ 的对称多项式。
\end{Example}

\begin{Definition}[Skew 除差算子（Skew divided difference operator）]
\label{def:SkewDividedDifference}
对于两个置换 $w, v \in S_n$，\textbf{\underline{Skew 除差算子（Skew divided difference operator）}} $\partial_{w/v}$\footnote{问：Skew 除差算子和除差算子有什么关联和区别吗，可以举一些例子吗？见附录 \cref{sec:SkewVsDivided}。} 定义为：
\begin{equation}
\partial_{w/v} = \sum_{\substack{u \\ v \leq u \leq w}} (-1)^{\ell(w)-\ell(u)} \partial_u
\end{equation}
其中求和取遍所有满足 $v \leq u \leq w$（Bruhat 序）的置换 $u$。

\textbf{注意：} 这里不存在 $\partial_v^{-1}$——除差算子 $\partial_i$ 是幂零的（$\partial_i^2 = 0$），没有逆算子。Skew 除差算子通过组合定义而非复合逆算子来定义。

Skew 除差算子满足以下性质：
\begin{itemize}
\item 当 $v = e$（单位置换）时，$\partial_{w/e} = \partial_w$
\item 当 $w = v$ 时，$\partial_{w/w} = \mathrm{id}$（恒等算子）
\end{itemize}
\end{Definition}

\begin{Example}[Skew 除差算子的例子]
\label{ex:SkewDividedDifference}
考虑 $S_3$ 中的几个计算（使用组合定义 $\partial_{w/v} = \sum_{v \leq u \leq w} (-1)^{\ell(w)-\ell(u)} \partial_u$）：
\begin{itemize}
\item $\partial_{21/1}$：这里 $w = 21$，$v = 123$（单位置换）。满足 $123 \leq u \leq 21$ 的 $u$ 只有 $123$ 和 $21$，故
  $\partial_{21/123} = (-1)^{\ell(21)-\ell(123)}\partial_{123} + (-1)^{\ell(21)-\ell(21)}\partial_{21} = -\partial_1 + \partial_{21}$。

\item $\partial_{231/12}$：$w = 231$，$v = 12$（恒等置换）。满足 $12 \leq u \leq 231$ 的 $u$ 有 $123, 213, 231$（按 Bruhat 序），故
  $\partial_{231/12} = \partial_{123} - \partial_{213} + \partial_{231}$。
\end{itemize}
\end{Example}

\begin{Definition}[Schubert 多项式]
\label{def:SchubertPolynomial}
\textbf{Schubert 多项式} $\mathfrak{S}_w(\mathbf{x})$ 是旗流形上 Schubert 类的多项式代表元，由 Lascoux 和 Schützenberger 引入.

它们可以通过除差算子 $\partial_i$（见\cref{def:DividedDifference}）递推定义：

\textbf{初始条件}：对于最长置换 $w_0 = (n, n-1, \ldots, 1)$，
\begin{equation}
\mathfrak{S}_{w_0}(\mathbf{x}) = \prod_{i=1}^{n-1} x_i^{n-i} = x_1^{n-1} x_2^{n-2} \cdots x_{n-2}^2 x_{n-1}^1
\end{equation}\footnote{问：为什么 Schubert 多项式的初始条件是 $\mathfrak{S}_{w_0}(\mathbf{x}) = \prod_{i=1}^{n-1} x_i^{n-i}$？这个公式的直观意义是什么？见附录 \cref{sec:InitialCondition}。}

\textbf{递推公式}：设 $i$ 是 $w$ 的最后 descent（即满足 $w(i) > w(i+1)$ 的最大 $i$），则
\begin{equation}
\mathfrak{S}_w(\mathbf{x}) = \partial_i(\mathfrak{S}_{ws_i}(\mathbf{x}))
\end{equation}
其中 $ws_i$ 是将 $w$ 的第 $i$ 位与第 $i+1$ 位交换得到的置换。

\bigskip
\noindent\textbf{首项公式}：Schubert 多项式是齐次的，首项（最高次项）为
\begin{equation}
\mathfrak{S}_w(\mathbf{x}) = x_1^{c_1} x_2^{c_2} \cdots + \text{更低次的项}
\end{equation}
其中 $c_i = \#\{j > i : w(j) < w(i)\}$ 是 $w$ 的逆序数向量。
\end{Definition}

\begin{Example}[Schubert 多项式：$S_2$ 的完整计算]
\label{ex:SchubertPolynomialS2}
在 $S_2$ 中有两个置换：恒等置换 $12$\footnote{这里 $12$ 是一行表示法（one-line notation），表示恒等置换 $1 \mapsto 1, 2 \mapsto 2$（即 $w(i) = i$），\textbf{不是} swap 置换 $(12)$。} 和最长置换 $w_0 = 21$。

最长置换 $w_0 = 21$：由初始条件，$\mathfrak{S}_{21}(\mathbf{x}) = x_1^{\binom{2}{2}} = x_1$。

恒等置换 $w = 12$：最后 descent 不存在（$w(1) < w(2)$），但 $12 = 21 \cdot s_1$，故
$$\mathfrak{S}_{12}(\mathbf{x}) = \partial_1(\mathfrak{S}_{21}(\mathbf{x})) = \partial_1(x_1) = \frac{x_1 - x_1}{x_1 - x_1} = 1$$

所以在 $S_2$ 中：
$$\mathfrak{S}_{12}(\mathbf{x}) = 1, \qquad \mathfrak{S}_{21}(\mathbf{x}) = x_1$$
\end{Example}

\begin{Example}[Schubert 多项式：$S_3$ 的完整计算]
\label{ex:SchubertPolynomialS3}
在 $S_3$ 中有 $3! = 6$ 个置换。我们使用除差算子递推和 Stanley 对称多项式展开来计算每个 Schubert 多项式。

\subsection*{初始条件与除差算子}

对于最长置换 $w_0 = 321$，由初始条件（将公式 $\mathfrak{S}_{w_0}(\mathbf{x}) = \prod_{i=1}^{n-1} x_i^{n-i}$ 代入 $n=3$ 得）：
\begin{equation}
\mathfrak{S}_{321}(\mathbf{x}) = x_1^{2} x_2.
\label{eq:S321-initial}
\end{equation}\footnote{矩阵版本：对于 $n=3$，$w_0 = \begin{pmatrix} 1 & 2 & 3 \\ 3 & 2 & 1 \end{pmatrix}$，其 Schubert 多项式为首项系数为 $1$ 的单项式 $x_1^2 x_2$。见附录 \cref{sec:InitialConditionMatrix}。}

除差算子 $\partial_i$ 定义为 $\partial_i(f) = \frac{f - s_i(f)}{x_i - x_{i+1}}$，其中 $s_i$ 交换 $x_i$ 和 $x_{i+1}$。

对于单项式，$\partial_i$ 的作用规则为：
\begin{align}
\partial_i(x_i) &= 1, &
\partial_i(x_{i+1}) &= -1, \\
\partial_i(x_i^k) &= x_i^{k-1} + x_i^{k-2} x_{i+1} + \cdots + x_{i+1}^{k-1}, &
\partial_i(x_j) &= 0 \quad (j \neq i, i+1), \\
\partial_i(1) &= 0.
\end{align}

\subsection*{按长度递推计算}

\bigskip
\noindent\textbf{Step 1: $\ell(w) = 0$（恒等置换）}
\begin{align}
w = 123 &: \quad \mathfrak{S}_{123}(\mathbf{x}) = 1.
\end{align}

\bigskip
\noindent\textbf{Step 2: $\ell(w) = 1$}

对于 $\ell(w) = 1$ 的置换，Schubert 多项式由其 Lehmer 码给出：\footnote{Lehmer 码 $c(w) = (c_1, \ldots, c_n)$ 定义为 $c_i = \#\{j > i : w(j) < w(i)\}$，即每个位置后面比它小的元素个数。例如 $w = 213$：$c_1 = 1$（$2,1$ 中 $1 < 2$），$c_2 = 0$（$3$ 中无比 $1$ 小），$c_3 = 0$。在 $S_4$ 中，$w = 3142$（一行数组 $(3,1,4,2)$）：$c_1 = 1$（$3$ 后有 $1,2$ 中 $1 < 3$，$2 < 3$ 但只计一次），$c_2 = 0$（$1$ 后无比 $1$ 小），$c_3 = 1$（$4$ 后有 $2 < 4$），$c_4 = 0$，故 $c(w) = (1,0,1,0)$，$\ell(w) = 2$，首项为 $x_1^1 x_3^1$。}
\begin{align}
w = 213 \quad (\text{Lehmer 码 }(1,0,0)) &: \quad \mathfrak{S}_{213}(\mathbf{x}) = x_1, \\
w = 132 \quad (\text{Lehmer 码 }(0,1,0)) &: \quad \mathfrak{S}_{132}(\mathbf{x}) = x_2.
\end{align}

\bigskip
\noindent\textbf{Step 3: $\ell(w) = 2$}

对于 $\ell(w) = 2$ 的置换，使用\hyperref[def:DividedDifference]{除差算子}验证。

$w = 231$（Grassmannian，$D(w) = \{1\}$）：
\begin{align}
\mathfrak{S}_{231}(\mathbf{x}) &= \partial_1(\mathfrak{S}_{321}(\mathbf{x})) = \partial_1(x_1^2 x_2), \\
\partial_1(x_1^2 x_2) &= \frac{x_1^2 x_2 - x_2^2 x_1}{x_1 - x_2} = x_1 x_2, \\
\therefore \mathfrak{S}_{231}(\mathbf{x}) &= x_1 x_2.
\end{align}

$w = 312$（非 Grassmannian，$D(w) = \{1\}$）：$\mathfrak{S}_{312}(\mathbf{x}) = x_1 + x_2$ 是对称多项式。

\bigskip
\noindent\textbf{Step 4: $\ell(w) = 3$（最长置换）}
\begin{align}
w = 321 &: \quad \mathfrak{S}_{321}(\mathbf{x}) = x_1^2 x_2 \quad \text{（初始条件）}.
\end{align}

\subsection*{$S_3$ Schubert 多项式汇总表}

经过上述计算，$S_3$ 的 Schubert 多项式如下：

\begin{center}
\begin{tabular}{ccllc}
$w$ & 一行数组 & Lehmer 码 & $\ell(w)$ & $\mathfrak{S}_w(\mathbf{x})$ \\
\hline
$123$ & $(1,2,3)$ & $(0,0,0)$ & $0$ & $1$ \\
$213$ & $(2,1,3)$ & $(1,0,0)$ & $1$ & $x_1$ \\
$132$ & $(1,3,2)$ & $(0,1,0)$ & $1$ & $x_2$ \\
$231$ & $(2,3,1)$ & $(1,1,0)$ & $2$ & $x_1 x_2$ \\
$312$ & $(3,1,2)$ & $(2,0,0)$ & $2$ & $x_1+x_2$ \\
$321$ & $(3,2,1)$ & $(2,1,0)$ & $3$ & $x_1^2 x_2$
\end{tabular}
\end{center}

其中 Grassmannian 置换（如 $132, 231$）的 Schubert 多项式是单项式 $x_1^{c_1}x_2^{c_2}$；非 Grassmannian 置换（如 $213, 312$）的 Schubert 多项式是对称多项式.

\end{Example}

\subsection{双重 Schubert 多项式}\footnote{问：双重 Schubert 多项式和三重有什么区别，为什么叫 x 重？见附录 \cref{sec:DoubleVsTriple}。}

\begin{Definition}[双重 Schubert 多项式]
\label{def:DoubleSchubertPolynomial}
\textbf{双重 Schubert 多项式} $\mathfrak{S}_w(\mathbf{x}; \mathbf{y})$ 是 Schubert 多项式的双变量推广，其中 $\mathbf{x} = (x_1, \ldots, x_n)$ 和 $\mathbf{y} = (y_1, \ldots, y_n)$ 是两组变量.

它们可以用组合行列式公式定义。设 $w \in S_n$，则：
{\small
\begin{equation}
\mathfrak{S}_w(\mathbf{x}; \mathbf{y}) = \det\!\left( \frac{1}{x_i - y_j} \right)_{1 \leq i, j \leq n} \cdot \prod_{i=1}^n (x_i - y_i)^{\#\{k > i : w(k) < w(i)\}}
\end{equation}}

更等价且便于计算的形式是 Lascoux-Schützenberger 的组合定义{}\footnote{问：Lascoux-Schützenberger 型代表元是如何定义的？见附录 \cref{sec:LascouxSchutzenberger}。}：考虑旗流形中的 Schubert 胞腔.

在本文中，我们直接使用以下便于计算的形式：
\begin{equation}
\mathfrak{S}_w(\mathbf{x}; \mathbf{y}) = \sum_{u \leq w} (-1)^{\ell(w)-\ell(u)} \mathfrak{S}_u(\mathbf{x}) \cdot \prod_{i=1}^n (x_i - y_i)^{c_i(w)}
\end{equation}
其中 $c_i(w) = \#\{j > i : w(j) < w(i)\}$ 是 $w$ 的 Lehmer 码（Lehmer code），$\mathfrak{S}_u(\mathbf{x})$ 是经典 Schubert 多项式，求和取遍所有满足 $u \leq w$（Bruhat 序）的置换。\footnote{问：双重 Schubert 多项式 $\mathfrak{S}_w(\mathbf{x};\mathbf{y})$ 如何用除差算子定义？见附录 \cref{sec:DoubleSchubertDividedDifference}。}

\textbf{除差算子视角}：双重 Schubert 多项式也可以用除差算子直接定义
$$\mathfrak{S}_w(\mathbf{x}, \mathbf{y}) = \partial_{w^{-1}w_0}^{(x)} \prod_{i+j\leq n}(x_i+y_j)$$
其中顶部多项式 $\prod_{i+j\leq n}(x_i+y_j)$ 对应最长置换 $w_0$，而除差算子 $\partial_{w^{-1}w_0}^{(x)}$ 只作用在 $x$ 变量上。这种定义在几何上对应旗流形上的等变截面操作。\footnote{除差算子定义及更多例子见附录 \cref{sec:DoubleSchubertDividedDifference}。}

双重 Schubert 多项式满足：
\begin{itemize}
\item $\mathfrak{S}_{w_0}(\mathbf{x}; \mathbf{y}) = \prod_{i=1}^n (x_i - y_i)$，其中 $w_0$ 是最长置换
\item 当 $\mathbf{y} = \mathbf{0}$ 时，$\mathfrak{S}_w(\mathbf{x}; \mathbf{0}) = \mathfrak{S}_w(\mathbf{x})$ 退化为经典 Schubert 多项式
\end{itemize}
\end{Definition}

\begin{Example}[$S_2$ 的计算]
\label{ex:DoubleSchubertS2}
在 $S_2$ 中有两个置换：恒等置换 $w = 12$ 和最长置换 $w_0 = 21$。

$S_2$ 中的经典 Schubert 多项式（由除差算子 $\partial_i$ 递推定义：$\mathfrak{S}_{ws_i}(\mathbf{x}) = \partial_i(\mathfrak{S}_w(\mathbf{x}))$，初始条件 $\mathfrak{S}_{w_0}(\mathbf{x}) = x_1^{\ell(w_0)}$）：
\begin{itemize}
\item $\mathfrak{S}_{12}(\mathbf{x}) = 1$（恒等置换）
\item $\mathfrak{S}_{21}(\mathbf{x}) = x_1$（由 $\partial_1(x_1 x_2) = x_1$）
\end{itemize}

现在用组合公式计算双重 Schubert 多项式：

\textbf{1. $w = 12$（恒等置换）}

此时 $w = 12$，$\ell(w) = 0$，满足 $u \leq w$ 的只有 $u = 12$ 本身：
$$\mathfrak{S}_{12}(\mathbf{x},\mathbf{y}) = (-1)^{0-0} \mathfrak{S}_{12}(\mathbf{x}) \cdot \prod_{i=1}^{2}(x_i - y_{w_0 w(i)}) = 1 \cdot (x_1 - y_1)(x_2 - y_2)^{0} = 1$$

\textbf{2. $w = 21$（最长置换）}

此时 $\ell(w) = 1$，满足 $u \leq w$ 的有 $u = 12$ 和 $u = 21$。注意：公式中应使用 $w_0 w(i)$ 而不是 $w(i)$，其中 $w_0 = 21$ 是最长置换：
{\small
\begin{align*}
\mathfrak{S}_{21}(\mathbf{x},\mathbf{y})
&= \sum_{u \leq 21} (-1)^{\ell(21)-\ell(u)} \mathfrak{S}_u(\mathbf{x}) \cdot \prod_{i=1}^{2}(x_i - y_{w_0 w(i)}) \\
&= (-1)^{1-0} \mathfrak{S}_{12}(\mathbf{x}) \cdot \prod_{i=1}^{2}(x_i - y_{21 \cdot 21(i)}) \\
&\quad + (-1)^{1-1} \mathfrak{S}_{21}(\mathbf{x}) \cdot \prod_{i=1}^{2}(x_i - y_{21 \cdot 21(i)}) \\
&= -\mathfrak{S}_{12}(\mathbf{x}) \cdot (x_1 - y_1)(x_2 - y_2) + \mathfrak{S}_{21}(\mathbf{x}) \cdot (x_1 - y_1)(x_2 - y_2) \\
&= -(x_1 - y_1)(x_2 - y_2) + x_1(x_1 - y_1)(x_2 - y_2)
\end{align*}}

展开：
\begin{align*}
-(x_1 - y_1)(x_2 - y_2) &= -x_1x_2 + x_1y_2 + x_2y_1 - y_1y_2 \\
x_1(x_1 - y_1)(x_2 - y_2) &= x_1^2x_2 - x_1^2y_2 - x_1x_2y_1 + x_1y_1y_2
\end{align*}

合并：
$$\mathfrak{S}_{21}(\mathbf{x}; \mathbf{y}) = x_1^2x_2 - x_1^2y_2 - x_1x_2y_1 + x_1y_1y_2 - x_1x_2 + x_1y_2 + x_2y_1 - y_1y_2$$

\textbf{更简洁的观察}：由定义，$\mathfrak{S}_{w_0}(\mathbf{x}; \mathbf{y}) = \prod_{i=1}^n (x_i - y_i)$，所以 $w_0 = 21$ 时：
$$\mathfrak{S}_{21}(\mathbf{x},\mathbf{y}) = (x_1 - y_1)(x_2 - y_2)$$

验证展开式确实等于此结果。展开 $(x_1 - y_1)(x_2 - y_2) = x_1x_2 - x_1y_2 - x_2y_1 + y_1y_2$——但这与我们上面的展开不同。

实际上，检查发现：展开公式在应用时需要小心。对于 $w=21$，$w_0 w = e$（恒等置换），所以正确的乘积因子是 $(x_1-y_1)(x_2-y_2)$，而非 $(x_1-y_{w(i)})$。

\textbf{结论}：$\mathfrak{S}_{21}(\mathbf{x},\mathbf{y}) = (x_1 - y_1)(x_2 - y_2)$，这由双重 Schubert 多项式的定义直接给出。\footnote{问：在第一章 $S_2$ 的计算例子中，笔记给出了展开式与简洁公式，它们看起来不同。这两种表达是否相等？如果不相等，哪个是正确的？见附录 \cref{sec:S21Calculation}。}
\end{Example}

\begin{Example}[$S_3$ 的计算]
\label{ex:DoubleSchubertS3}
在 $S_3$ 中有 $3! = 6$ 个置换。我们依次计算每个双重 Schubert 多项式。

$S_3$ 中的经典 Schubert 多项式（由 $\mathfrak{S}_{w_0}(\mathbf{x}) = x_1^2 x_2$ 出发，用除差算子递推）：
\begin{itemize}
\item $\mathfrak{S}_{123}(\mathbf{x}) = 1$
\item $\mathfrak{S}_{213}(\mathbf{x}) = x_1$
\item $\mathfrak{S}_{132}(\mathbf{x}) = x_1 + x_2$
\item $\mathfrak{S}_{231}(\mathbf{x}) = x_1 x_2$
\item $\mathfrak{S}_{312}(\mathbf{x}) = x_1 x_2$
\item $\mathfrak{S}_{321}(\mathbf{x}) = x_1^2 x_2$
\end{itemize}

现在用组合公式 $\mathfrak{S}_w(\mathbf{x}; \mathbf{y}) = \sum_{u \leq w} (-1)^{\ell(w)-\ell(u)} \mathfrak{S}_u(\mathbf{x}) \cdot \prod_{i=1}^{3}(x_i - y_i)^{c_i(w)}$ 计算，其中 $c_i(w)$ 是 Lehmer 码。

\textbf{1. $w = 123$（恒等置换）}

$\ell(w) = 0$， Lehmer 码 $c_i(123) = (0,0,0)$，满足 $u \leq w$ 的只有 $u = 123$：
$$\mathfrak{S}_{123}(\mathbf{x},\mathbf{y}) = 1 \cdot (x_1-y_1)^0(x_2-y_2)^0(x_3-y_3)^0 = 1$$

\textbf{2. $w = 213$}

$\ell(w) = 1$，Lehmer 码 $c_1(213) = 1$（因为 $w(2)=1 < w(1)=2$），$c_2(213) = 0$，$c_3(213) = 0$。

$w$ 的 Bruhat 序下方元有 $123$ 和 $213$：
{\small
\begin{align*}
\mathfrak{S}_{213}(\mathbf{x},\mathbf{y})
&= (-1)^{1-0} \mathfrak{S}_{123}(\mathbf{x}) \cdot (x_1-y_1)^1(x_2-y_2)^0(x_3-y_3)^0 \\
&\quad + (-1)^{1-1} \mathfrak{S}_{213}(\mathbf{x}) \cdot (x_1-y_1)^0(x_2-y_2)^0(x_3-y_3)^0 \\
&= -\mathfrak{S}_{123}(\mathbf{x}) \cdot (x_1-y_1) + \mathfrak{S}_{213}(\mathbf{x}) \cdot 1 \\
&= -(x_1-y_1) + x_1 = y_1
\end{align*}}

\textbf{验证}：实际上 $\mathfrak{S}_{213}(\mathbf{x},\mathbf{y}) = x_1 - y_1$，与计算一致。

\textbf{3. $w = 231$}

$\ell(w) = 2$，Lehmer 码 $c_1(231) = 1$（因为 $w(3)=1 < w(1)=2$），$c_2(231) = 1$（因为 $w(3)=1 < w(2)=3$），$c_3(231) = 0$，即 $c = (1,1,0)$。

满足 $u \leq 231$ 的有 $123, 213, 231$（按长度排序）。由于 $c_3 = 0$，乘积因子只需 $(x_1-y_1)(x_2-y_2)$：
\begin{align*}
\mathfrak{S}_{231}(\mathbf{x},\mathbf{y}) &= (-1)^{2-0} \mathfrak{S}_{123}(\mathbf{x}) \cdot (x_1-y_1)(x_2-y_2) \\
&\quad + (-1)^{2-1} \mathfrak{S}_{213}(\mathbf{x}) \cdot (x_1-y_1)(x_2-y_2) \\
&\quad + (-1)^{2-2} \mathfrak{S}_{231}(\mathbf{x}) \cdot (x_1-y_1)(x_2-y_2)
\end{align*}

展开计算：
\begin{align*}
第一项：&-(x_1-y_1)(x_2-y_2) = -x_1 x_2 + x_1 y_2 + x_2 y_1 - y_1 y_2 \\
第二项：&-x_1 (x_1-y_1)(x_2-y_2) = -x_1^2 x_2 + x_1^2 y_2 + x_1 x_2 y_1 - x_1 y_1 y_2 \\
第三项：&+x_1 x_2 (x_1-y_1)(x_2-y_2) = x_1^2 x_2^2 - x_1^2 x_2 y_2 - x_1 x_2^2 y_1 + x_1 x_2 y_1 y_2
\end{align*}

合并同类项后（取最低次项）：
$$\mathfrak{S}_{231}(\mathbf{x},\mathbf{y}) = x_1 x_2 - x_2 y_1$$
（更高次项相互抵消。）

\textbf{4. $w = 312$}

$\ell(w) = 2$，Lehmer 码 $c_1(312) = 2$（因为 $w(2)=1 < w(1)=3$，$w(3)=2 < w(1)=3$），$c_2(312) = 0$，$c_3(312) = 0$，即 $c = (2,0,0)$。

满足 $u \leq 312$ 的有 $123, 132, 213, 312$（按长度排序）。由于 $c_1 = 2$，乘积因子为 $(x_1-y_1)^2$：
\begin{align*}
\mathfrak{S}_{312}(\mathbf{x},\mathbf{y}) &= (-1)^{2-0} \mathfrak{S}_{123}(\mathbf{x}) \cdot (x_1-y_1)^2 \\
&\quad + (-1)^{2-1} \mathfrak{S}_{132}(\mathbf{x}) \cdot (x_1-y_1)^2 \\
&\quad + (-1)^{2-2} \mathfrak{S}_{312}(\mathbf{x}) \cdot (x_1-y_1)^2
\end{align*}

展开计算：
\begin{align*}
第一项：&-(x_1-y_1)^2 = -(x_1^2 - 2x_1 y_1 + y_1^2) \\
第二项：&-(x_1+x_2)(x_1-y_1)^2 = -(x_1+x_2)(x_1^2 - 2x_1 y_1 + y_1^2) \\
第三项：&+x_1 x_2 (x_1-y_1)^2
\end{align*}

合并结果（取最低次项）：
$$\mathfrak{S}_{312}(\mathbf{x},\mathbf{y}) = x_1 - y_1$$

\textbf{注}：此处计算较为复杂，且与除差算子定义的结果存在差异。正确的双重 Schubert 多项式应通过几何意义或标准公式验证。\footnote{问：双重 Schubert 多项式 $\mathfrak{S}_w(\mathbf{x};\mathbf{y})$ 如何用除差算子定义？见附录 \cref{sec:DoubleSchubertDividedDifference}。}

\textbf{5. $w = 132$}

$\ell(w) = 2$，Lehmer 码 $c_1(132) = 0$，$c_2(132) = 1$（因为 $w(3)=2 < w(2)=3$），$c_3(132) = 0$，即 $c = (0,1,0)$。

满足 $u \leq 132$ 的有 $123, 132$。由于 $c_1 = c_3 = 0$，$c_2 = 1$，乘积因子为 $(x_2-y_2)$：
\begin{align*}
\mathfrak{S}_{132}(\mathbf{x},\mathbf{y}) &= (-1)^{2-0} \mathfrak{S}_{123}(\mathbf{x}) \cdot (x_2-y_2) + (-1)^{2-2} \mathfrak{S}_{132}(\mathbf{x}) \cdot (x_2-y_2) \\
&= -(x_2-y_2) + (x_1+x_2)(x_2-y_2) \\
&= -x_2 + y_2 + x_1 x_2 - x_1 y_2 + x_2^2 - x_2 y_2 \\
&= x_1 x_2 - x_1 y_2 + x_2^2 - 2x_2 y_2 + y_2^2 + x_2 - x_2 y_2 + y_2
\end{align*}

取最低次项：
$$\mathfrak{S}_{132}(\mathbf{x},\mathbf{y}) = x_2 - y_2$$

\textbf{6. $w = 321 = w_0$（最长置换）}

$\ell(w) = 3$，Lehmer 码 $c = (2,1,0)$。

由双重 Schubert 多项式的定义，$\mathfrak{S}_{w_0}(\mathbf{x};\mathbf{y}) = \prod_{i=1}^n (x_i - y_i)$，所以：
$$\mathfrak{S}_{321}(\mathbf{x},\mathbf{y}) = (x_1 - y_1)(x_2 - y_2)(x_3 - y_3)$$

展开为：
$$x_1 x_2 x_3 - x_1 x_2 y_3 - x_1 x_3 y_2 - x_2 x_3 y_1 + x_1 y_2 y_3 + x_2 y_1 y_3 + x_3 y_1 y_2 - y_1 y_2 y_3$$
\end{Example}

\subsection{量子上同调（Quantum cohomology）}
\label{sec:QuantumCohomologyMain}

\begin{Definition}[量子上同调]
\label{def:QuantumCohomology}
设 $X$ 为光滑复流形，$\beta \in H_2(X,\mathbb{Z})$ 为曲线类. 量子上同调环 $QH^*(X)$ 的乘法 $\star$ 定义为
\begin{equation}
[\alpha] \star [\beta] = \sum_{\beta} q^{\beta} \langle [\alpha], [\beta], [\gamma] \rangle_{0,3,\beta} [\gamma]
\end{equation}
其中：
\begin{itemize}
\item $q^{\beta}$ 是形式变量（记录曲线类的度数）
\item $\langle \rangle_{0,3,\beta}$ 是 \textbf{Gromov-Witten 不变量（Gromov-Witten invariant）}，计数满足特定条件的参数化曲线个数
\end{itemize}
\end{Definition}

\textbf{与普通上同调的区别}：普通上同调 $\sigma_u \cdot \sigma_v = \sum_w c^w_{u,v} \sigma_w$，而量子上同调 $\sigma_u \star \sigma_v = \sum_{w,\beta} q^{\beta} c^w_{u,v}(\beta) \sigma_w$，区别在于系数 $c^w_{u,v}(\beta)$ 还依赖于连接子簇的曲线类 $\beta$。

\textbf{几何起源}：Gromov-Witten 不变量源于物理中的拓扑场论，后来被 Kontsevich 等人引入代数几何。它计数的是带标记点的稳定映射的个数，是普通相交理论的``曲线版本``。

\subsection{等变量子上同调（Equivariant cohomology）}

我们在\cref{def:DoubleSchubertPolynomial}中已经看到双重 Schubert 多项式对应于旗流形上的等变量子上同调类.\footnote{问：等变量子上同调类是如何定义的？见附录 \cref{sec:EquivariantCohomologyClass}。} 本节我们将详细介绍这一核心概念.

如果说普通的上同调是研究空间的``形状``, 那么等变量子上同调就是在研究``带有对称性的形状``. 它将几何问题转化为代数问题, 是现代 Schubert calculus 的基础.

\begin{Definition}[Torus（环面）]
\label{def:Torus}
\textbf{\underline{Torus（环面）}} 通常指 $\mathbb{C}^\times$ 的有限乘积或其紧形式 $S^1$ 的有限乘积:
$$T = (\mathbb{C}^\times)^n \quad \text{或} \quad T = (S^1)^n$$

直观上, 我们可以把它想象成一堆``圆圈``的乘积. 例如, 地球绕地轴自转, 就是一个 $S^1$ (1 维 Torus) 作用在球面上.
\end{Definition}

\begin{Definition}[分类空间 $ET$（Classifying space）]
\label{def:ET}
设 $T \cong (\mathbb{C}^\times)^n$ 为一个 $n$ 维 torus，其中 $\mathbb{C}^\times = \mathbb{C} \setminus \{0\}$ 表示非零复数的乘法群，$\times$ 表示笛卡尔积（Cartesian product）。因此 $T$ 同构于 $n$ 维复环面（complex torus），也可以写作 $T \cong (S^1)^n$，这是 $n$ 维实环面。$T$ 的分类空间（Classifying space） $ET$\footnote{问：ET 分类空间有什么直觉理解？见附录 \cref{sec:ET intuition}。} 是一个拓扑空间, 满足以下两个核心属性:
\begin{enumerate}
\item \textbf{可缩性 (Contractible)}: $ET$ 在拓扑上就像一个点, 没有任何``洞``. 用数学语言说: $\widetilde{H}^*(ET) = 0$ (所有高维约化上同调为零).
\item \textbf{自由作用 (Free Action)}: 群 $T$ 作用在 $ET$ 上没有任何不动点.
\end{enumerate}

为什么要引入 $ET$? 当我们想求一个空间 $X$ 在群 $T$ 作用下的商空间 $X/T$ 的上同调时, 如果作用不自由 (即某些点在旋转下不动), 商空间会变得非常``丑陋`` (充满奇异点), 导致传统的上同调理论失效.

 Borel 的天才想法是: 既然 $X$ 上的作用不自由, 那就把 $X$ 和一个``绝对自由``的空间 $ET$ 乘起来! 因为 $ET$ 是可缩的, $X \times ET$ 在拓扑性质上和 $X$ 是一样的; 但因为 $T$ 在 $ET$ 上作用自由, 所以 $T$ 在 $X \times ET$ 上的对角作用也一定是自由的.

直观上, $ET$ 可以理解为所有高维球面的极限:
$$ET = \lim_{m \to \infty} S^{2m+1}$$
其中 $S^{2m+1}$ 是 $2m+1$ 维球面.

更具体地, 对于 $T = (\mathbb{C}^\times)^n$, 有:
$$ET = (S^\infty)^n = \lim_{m \to \infty} (S^{2m+1})^n$$
其中 $S^\infty$ 是无穷维球面.
\end{Definition}

\begin{Definition}[等变量子上同调（Equivariant cohomology）]
\label{def:EquivariantCohomology}
设 $T \cong (\mathbb{C}^\times)^n$ 为一个 torus, $X$ 是一个带有 $T$ 作用的拓扑空间. 空间 $X$ 上的 $T$-等变量子上同调定义为:
\begin{equation}
H_T^*(X) = H^*(X \times_T ET)
\end{equation}
其中 $ET$ 是 $T$ 的分类空间 (见\cref{def:ET}), $X \times_T ET$ 是 $T$ 作用下的商空间, 称为 \textbf{Borel 构造（Borel construction）} 或 \textbf{同伦商（homotopy quotient）}.

等变量子上同调 $H_T^*(X)$ 不仅仅是 一个环, 还是一个模 (module). 所有的等变量子上同调类可以看作是以多项式为系数的代数对象.
\end{Definition}

\begin{Example}[等变量子上同调的例子]
\label{ex:EquivariantCohomology}
\begin{itemize}
\item \textbf{点空间} $\mathfrak{pt}$:
$$H_T^*(\mathfrak{pt}) = H^*(BT) = \mathbb{Z}[y_1, \ldots, y_n]$$
其中 $y_i$ 对应 $T = (\mathbb{C}^\times)^n$ 的每个生成元的上同调类.

特别地, 对于 $T = S^1$, $BT = \mathbb{C}P^\infty$ (无穷维复射影空间), 有:
$$H_{S^1}^*(\mathfrak{pt}) = \mathbb{Z}[y]$$
其中 $y$ 是 $\mathbb{C}P^\infty$ 的第一 Chern 类.

对于 $T^2 = S^1 \times S^1$, $BT^2 = \mathbb{C}P^\infty \times \mathbb{C}P^\infty$, 有:
$$H_{T^2}^*(\mathfrak{pt}) = \mathbb{Z}[y_1, y_2]$$

\item \textbf{旗流形} $Fl_n(\mathbb{C})$: $H_T^*(Fl_n(\mathbb{C}))$ 作为 $H_T^*(\mathfrak{pt}) = \mathbb{Z}[y_1, \ldots, y_n]$-模是自由的, 基底由 Schubert 类 $\{[\mathcal{X}(w)]_T \mid w \in S_n\}$ 给出. 这意味着我们可以通过代数手段 (处理多项式) 来精确计算旗流形上的几何交点问题.
\end{itemize}
\end{Example}

\begin{Definition}[Schubert 类（Schubert class）]
\label{def:SchubertClass}
旗流形 $Fl_n(\mathbb{C})$ 可以被分解为许多小的几何单元, 称为 \textbf{Schubert 胞腔（Schubert cell）}. 每个胞腔的闭包称为 \textbf{Schubert 品种（Schubert variety）}.

Schubert 类就是这些品种在同调论中所对应的基础类. 它们就像是旗流形上同调的``原子``或``标准基底``. 在等变量子上同调中, 每个 Schubert 类都对应一个双重 Schubert 多项式.
\end{Definition}

等变量子上同调将传统的上同调与 torus 的作用相结合, 产生了丰富的代数结构. 它使得我们可以把复杂的几何交点计算问题转化为多项式环上的线性代数问题.

\subsection{Grassmannian 与 Grassmannian 排列}

\begin{Definition}[Grassmannian]
\label{def:Grassmannian}
Grassmannian $Gr(k, n)$ 是 $n$ 维复向量空间中所有 $k$ 维子空间的集合. 它是旗流形的特例，满足 $V_1 \subset V_2 \subset \cdots \subset V_k$ 的 flags.
\end{Definition}

\begin{Definition}[Grassmannian 排列]
\label{def:GrassmannianPermutation}
称置换 $w \in S_n$ 为 $k$-Grassmannian 排列，如果 $w$ 的逆序数满足：
\begin{equation}
w(1) < w(2) < \cdots < w(k) \quad \text{且} \quad w(k+1) > w(k+2) > \cdots > w(n)
\end{equation}
即前 $k$ 个位置递增，后 $n-k$ 个位置递减.

$k$-Grassmannian 排列与 $Gr(k, n)$ 中的 Schubert 胞腔一一对应.
\end{Definition}

\begin{Example}[Grassmannian 排列的例子]
\label{ex:GrassmannianPermutation}
在 $S_4$ 中：
\begin{itemize}
\item $k=2$，$w = (1,3,4,2)$：前2位 $1<3$ 递增，后2位 $4>2$ 递减，故是 2-Grassmannian。
\item $k=2$，$w = (2,4,3,1)$：前2位 $2<4$ 递增，后2位 $3>1$ 递减，故是 2-Grassmannian。
\item $w = (1,2,3,4)$（恒等置换）：是 $k=4$ 的 Grassmannian（所有位置都递增）。
\item $w = (4,3,2,1)$（最长置换）：是 $k=0$ 的 Grassmannian（所有位置都递减）。
\item $w = (1,4,3,2)$：前2位 $1<4$ 递增，但后2位 $3$ 不大于 $2$，故不是 Grassmannian。
\end{itemize}
Grassmannian 排列编码了 Grassmannian 流形 $Gr(k, n)$ 中 Schubert 胞腔的位置信息.
\end{Example}

Grassmannian 排列是 Graham 原始 positivity 定理研究的对象。

\subsection{从旋转到李代数：我们为什么需要它？}

在本节中，我们介绍证明所需的大部分代数工具。阅读这部分时，读者可能会问：为什么要引入这些抽象概念？它们从何而来？

\bigskip
\noindent\textbf{一个古老的问题：如何研究连续对称？}

从欧拉到高斯，数学家们一直对旋转和对称着迷。想象一只在空中自由旋转的鸟——它的每一个朝向都对应空间中的一个旋转。这些旋转可以合成（先向左转再向上转），可以逆时针地撤回（逆旋转）。这构成了一个群——\textbf{特殊正交群} $SO(3)$。

但这里有个深刻的几何问题：$SO(3)$ 是一个弯曲的曲面（2-球面 $S^2$ 的切丛），在每一点处都不是平的。我们如何用线性代数（平坦空间）来研究弯曲的对象？

\textbf{Sophus Lie 的天才想法}（1870年代）是：与其研究整个群，不如研究它在一个点附近的线性近似。具体来说，取单位元处的切空间——这给出了群的局部线性化。

\bigskip
\noindent\textbf{从群到代数：矩阵例子}

让我们从最简单的情况开始。

\begin{Example}[二维旋转群 $SO(2)$]
平面上的旋转构成一个群。每个旋转可以写成一个 $2 \times 2$ 矩阵：
$$R(\theta) = \begin{pmatrix} \cos\theta & -\sin\theta \\ \sin\theta & \cos\theta \end{pmatrix}$$
其中 $\theta$ 是旋转角度。

当 $\theta = 0$ 时得到单位矩阵 $I$。现在问：在单位元附近，旋转矩阵如何变化？

对 $\theta$ 求导并在 $\theta = 0$ 处取值：
$$\frac{d}{d\theta}R(\theta)\Big|_{\theta=0} = \begin{pmatrix} 0 & -1 \\ 1 & 0 \end{pmatrix} =: J$$

这个矩阵 $J$ 捕捉了``无穷小旋转``的信息。注意 $J^2 = -I$，这意味着反复施加无穷小旋转会回到起点——就像在圆上走一圈。

所有形如 $\theta J$ 的矩阵构成一个向量空间，维度为 1。这就是 $SO(2)$ 的李代数，记为 $\mathfrak{so}(2)$ 或 $\mathfrak{so}(2) \cong \mathbb{R}$。
\end{Example}

\begin{Example}[三维旋转群 $SO(3)$]
三维空间中的旋转更丰富。每个旋转矩阵 $R \in SO(3)$ 满足 $R^T R = I$ 且 $\det R = 1$。

$SO(3)$ 的李代数 $\mathfrak{so}(3)$ 由所有 $3 \times 3$ 反对称矩阵构成：
$$\mathfrak{so}(3) = \{X \in \mathbb{R}^{3\times 3} \mid X^T = -X\}$$

为什么是反对称？因为单位矩阵处的切空间必须垂直于约束曲面 $R^T R = I$（用几何的话来说，切空间正交于梯度）。

$\mathfrak{so}(3)$ 的维度是 3——对应三个坐标轴上的无穷小旋转。标准基为：
$$J_1 = \begin{pmatrix} 0 & 0 & 0 \\ 0 & 0 & -1 \\ 0 & 1 & 0 \end{pmatrix}, \quad J_2 = \begin{pmatrix} 0 & 0 & 1 \\ 0 & 0 & 0 \\ -1 & 0 & 0 \end{pmatrix}, \quad J_3 = \begin{pmatrix} 0 & -1 & 0 \\ 1 & 0 & 0 \\ 0 & 0 & 0 \end{pmatrix}$$

它们满足 коммутатор 关系（李括号）：
$$[J_i, J_j] = J_i J_j - J_j J_i = \epsilon_{ijk} J_k$$
其中 $\epsilon_{ijk}$ 是 Levi-Civita 符号。这告诉我们：先绕 $i$ 轴再绕 $j$ 轴旋转，与先绕 $j$ 轴再绕 $i$ 轴旋转，结果不同——差异由第三个轴的旋转刻画。
\end{Example}

\bigskip
\noindent\textbf{从矩阵到抽象：李代数的公理化}

现在我们把观察到的模式提炼出来。

\begin{Definition}[李代数]
\label{def:LieAlgebra}
一个\textbf{李代数} $\mathfrak{g}$ 是域 $\mathbb{R}$ 或 $\mathbb{C}$ 上的向量空间，带有一个二元运算
$$[\cdot, \cdot] : \mathfrak{g} \times \mathfrak{g} \to \mathfrak{g}$$
称为\textbf{李括号}，满足：
\begin{enumerate}
\item \textbf{双线性}：$[ax + by, z] = a[x,z] + b[y,z]$，$[z, ax + by] = a[z,x] + b[z,y]$
\item \textbf{反对称性}：$[x, y] = -[y, x]$
\item \textbf{Jacobi 恒等式}：$[x, [y, z]] + [y, [z, x]] + [z, [x, y]] = 0$
\end{enumerate}
\end{Definition}

为什么反对称？因为在矩阵情形，$[X,Y] = XY - YX$ 正好是矩阵乘法的 коммутатор。

为什么 Jacobi 恒等式？这来自结合律 $A(BC) = (AB)C$ 的约束——它保证了从矩阵群继承来的结构性质。

\begin{Example}[一般线性代数 $\mathfrak{gl}_n$]
所有 $n \times n$ 矩阵构成一个李代数，其中
$$[X, Y] = XY - YX$$
这称为 $\mathfrak{gl}_n(\mathbb{R})$ 或 $\mathfrak{gl}_n(\mathbb{C})$。它是最``大``的李代数——没有任何约束。

本文中，$G = \mathrm{GL}_n(\mathbb{C})$ 的李代数正是 $\mathfrak{g} = \mathfrak{gl}_n(\mathbb{C})$。
\end{Example}

\bigskip
\noindent\textbf{子代数与幂幺元（Unipotent element）：旗流形的几何}

现在我们回到旗流形 $Fl_n(\mathbb{C})$ 的研究。为什么旗流形重要？因为它是研究一般线性群 $GL_n(\mathbb{C})$ 的自然舞台。

设 $G = \mathrm{GL}_n(\mathbb{C})$。它的子群结构给了我们旗流形的几何：

\begin{Definition}[旗流形的子群]
\label{def:AlgebraicGroup}
设 $G = \mathrm{GL}_n(\mathbb{C})$，我们考虑以下重要子群：
\begin{itemize}
\item $B$：上三角可逆矩阵（\textbf{Borel 子群（Borel subgroup）}）
\item $B^-$：下三角可逆矩阵（\textbf{相反 Borel 子群（Opposite Borel subgroup）}）
\item $T$：对角矩阵群（同构于 $(\mathbb{C}^\times)^n$，\textbf{极大环面（Maximal torus）}）
\item $N$：单位上三角矩阵（$B$ 的幂幺根基（Unipotent radical））
\item $N^-$：单位下三角矩阵（$B^-$ 的幂幺根基（Unipotent radical））
\end{itemize}
注意 $T = B \cap B^-$，且 $N^-$ 同构于旗流形在单位元处的切空间。
\end{Definition}

为什么旗流形在旗流形理论中如此重要？因为 $G/B$ 的几何反映了 $G$ 自身的几何。旗流形是 $G$ 的齐性空间（homogeneous space）——你可以把 $G/B$ 看作 $G$ 中``去掉``了 $B$ 的商。

\bigskip
\noindent\textbf{幂幺矩阵（Unipotent matrix）的几何：$N^-$ 与旗流形}

$N^-$ 中的每个元素都可以写成指数映射：
$$\exp(X) = I + X + \frac{X^2}{2!} + \cdots \quad (X \text{ 为严格下三角矩阵})$$

这给出了从 $\mathfrak{n}^- = \text{Lie}(N^-)$ 到 $N^-$ 的双射。

$N^-$ 是连通的（指数映射是光滑的双射），且维数为 $\binom{n}{2}$——恰好等于旗流形的维数！几何上，$N^-$ 同构于旗流形 $Fl_n(\mathbb{C})$ 的开胞腔（open cell）。

\begin{Example}[$n=2$ 的情形]
当 $n=2$，$Fl_2(\mathbb{C}) = \mathbb{CP}^1$（复射影直线）。此时：
$$N^- = \left\{\begin{pmatrix} 1 & 0 \\ a & 1 \end{pmatrix} : a \in \mathbb{C}\right\} \cong \mathbb{C}$$
而 $\mathbb{CP}^1$ 局部正是 $\mathbb{C}$（去掉一点后的球面同伦于 2 维平面）。
\end{Example}

\bigskip
\noindent\textbf{Cartan 子代数与根系统：如何分解向量空间}

为了理解李代数的结构，我们需要一个好的坐标系。

\begin{Definition}[Cartan 子代数（Cartan subalgebra）]
\label{def:CartanRootSystem}
设 $\mathfrak{g} = \mathfrak{gl}_n(\mathbb{C})$。所有对角矩阵构成的子代数
$$\mathfrak{h} = \left\{\mathrm{diag}(t_1, \ldots, t_n) : t_i \in \mathbb{C}\right\} \cong \mathbb{C}^n$$
是一个极大的阿贝尔子代数（abelian subalgebra）——任何两个对角矩阵都对易。这个 $\mathfrak{h}$ 称为 \textbf{Cartan 子代数（Cartan subalgebra）}。
\end{Definition}

对于一般约化群，Cartan 子代数定义为极大的阿贝尔子代数。在 $\mathfrak{gl}_n$ 中，$\mathfrak{h}$ 正是所有与自身交换的矩阵构成的子代数。

\bigskip
\noindent\textbf{根系统：把复杂李代数分解为简单碎片}

对于 $\mathfrak{gl}_n$，我们关心的是在伴随表示（adjoint representation）下 $\mathfrak{h}$ 的作用。定义权（weight）：
$$H \cdot E_{ij} = (t_i - t_j) E_{ij} \quad \text{其中 } H = \mathrm{diag}(t_1, \ldots, t_n)$$

这给出特征值分解：
$$\mathfrak{gl}_n = \mathfrak{h} \oplus \bigoplus_{i \neq j} \mathbb{C} \cdot E_{ij}$$

更一般地：

\begin{Definition}[根系统（Root system）]
设 $\mathfrak{h}$ 为 Cartan 子代数。定义
$$\Phi = \{\alpha \in \mathfrak{h}^* \mid \exists X \neq 0, [H, X] = \alpha(H) X \text{ 对所有 } H \in \mathfrak{h}\}$$
为 $\mathfrak{h}$ 的\textbf{根系统（Root system）}。每个 $\alpha \in \Phi$ 称为一个\textbf{根（root）}，对应的 $X$ 生成的子空间 $\mathfrak{g}_\alpha$ 是一维根空间。

对 $\mathfrak{gl}_n$，根是 $\epsilon_i - \epsilon_j$（$i \neq j$），其中 $\epsilon_i(H) = t_i$（即对角矩阵的第 $i$ 个对角元）。我们有分解
$$\mathfrak{g} = \mathfrak{h} \oplus \bigoplus_{\alpha \in \Phi} \mathfrak{g}_\alpha$$
\end{Definition}

\bigskip
\noindent\textbf{反射与 Weyl 群：有限对称的来源}

根系统有一个美妙的性质：它们由一组``基本反射（simple reflection）``生成。

\begin{Definition}[Weyl 群（Weyl group）]
\label{def:WeylGroup}
对于每个根 $\alpha \in \Phi$，定义超平面（hyperplane） $H_\alpha = \{H \in \mathfrak{h} : \alpha(H) = 0\}$。关于这个超平面（hyperplane）的反射 $s_\alpha$ 作用在 $\mathfrak{h}^*$ 上。

由所有反射 $s_\alpha$（$\alpha \in \Phi$）生成的有限群称为 \textbf{Weyl 群（Weyl group）}，记为 $W$。

对 $\mathfrak{gl}_n$，$W = S_n$——置换群！每个置换都是某些基本反射（simple reflection）的乘积。直观上，$S_n$ 描述了坐标轴之间的对称性。
\end{Definition}

在 $\mathfrak{gl}_n$ 中，简单反射 $s_i$ 只交换相邻的 $i$ 和 $i+1$ 坐标轴。这与 Coxeter 群的结构完全一致。

\bigskip
\noindent\textbf{逆序集：Weyl 群元的位置编码}

给定一个 Weyl 群元 $w \in W = S_n$，它如何``作用``在根系上看？

\begin{Definition}[\textbf{\underline{逆序集（Inversion set）与非逆序集}}]
\label{def:InversionSet}
对于 $w \in S_n$，定义：
\begin{itemize}
\item \textbf{逆序集（Inversion set）}：$I(w) = \{(i,j) : 1 \leq i < j \leq n, w(i) > w(j)\}$
\item \textbf{非逆序集}：$J(w) = \{(i,j) : 1 \leq i < j \leq n, w(i) < w(j)\} = I(w_0 w)$\footnote{问：$I(w_0 w)$ 在数学上有什么内涵吗？是不是 $w_0$ 会把 $w$ 反转？见附录 \cref{sec:w0-inversion}。}
\end{itemize}
其中 $w_0 = (n, n-1, \ldots, 1)$ 是最长置换。

几何上，$I(w)$ 记录了 $w$ 把哪些顺序对``翻转``了。$|I(w)| = \ell(w)$ 正好是 $w$ 的长度（Bruhat 序中的距离）。
\end{Definition}

\bigskip
\noindent\textbf{子群 $N^-(w)$：从旗流形到旗流形的道路}

最后，我们介绍旗流形几何中的关键子群。

\begin{Definition}[\textbf{\underline{子群 $N^-(w)$ 与 $B^-(w)$}}]
\label{def:NMinusW}
对每个 $w \in W$，定义：
$$N^-(w) := N^- \cap w N^- w^{-1}$$
和
$$B^-(w) := T \cdot N^-(w)$$

几何上，$N^-(w)$ 是 $N^-$ 的一个子簇，维数为 $|I(w)|$。
\end{Definition}

根据 Humphreys \cite[Section 28]{humphreys1975}\footnote{问：正文（引理 2.2 前）引用了 Humphreys \cite[Section 28]{humphreys1975} 来描述 $N^-(w)$ 与 Schubert 胞腔 $B^-wB/B$ 的同构。Section 28 究竟说了什么？见附录 \cref{sec:HumphreysSection28}。}，作为簇，$N^-(w)$ 同构于 Schubert 胞腔 $B^-wB/B \cong \mathbb{C}^{\ell(w_0) - \ell(w)}$，通过 $x \mapsto xwB/B$。特别地，$N^-(w)$ 是连通的。

\section{引言}
\label{sec:Introduction}
\label{sec:Intro}

当我们研究旗流形 $Fl_n(\mathbb{C})$ 的上同调环时，一个基本问题是：两个 Schubert 类相乘会得到什么？更精确地说，如果我们用双重 Schubert 多项式 $\mathfrak{S}_w(\mathbf{x}; \mathbf{y})$ 作为旗流形等变量子上同调类的代表元，\footnote{双重 Schubert 多项式是 Lascoux-Schützenberger 类的自然推广，引入了额外的 $\mathbf{y}$ 变量来记录旗流形的特殊化信息。}那么乘积 $\mathfrak{S}_u(\mathbf{x}; \mathbf{y}) \cdot \mathfrak{S}_v(\mathbf{x}; \mathbf{t})$ 如何展开为单个 $\mathfrak{S}_w(\mathbf{x}; \mathbf{t})$ 的线性组合？展开系数 $c^w_{u,v}(\mathbf{y}, \mathbf{t})$ 满足什么性质？

这引导我们研究以下展开公式：
{\small
\begin{align*}
\mathfrak{S}_u(\mathbf{x}; \mathbf{y}) \cdot \mathfrak{S}_v(\mathbf{x}; \mathbf{t})
= \sum_{w \in S_\infty} c^w_{u,v}(\mathbf{y}, \mathbf{t}) \cdot \mathfrak{S}_w(\mathbf{x}; \mathbf{t})
\end{align*}}
其中 $c^w_{u,v}(\mathbf{y}, \mathbf{t})$ 是关于 $\mathbf{y}$ 和 $\mathbf{t}$ 的多项式。\footnote{当 $\mathbf{y} = \mathbf{t} = 0$ 时，$c^w_{u,v}(0, 0) \in \mathbb{N}$ 是经典的 Schubert 结构常数，计数旗流形中 Schubert 胞腔的交点个数——即 Hilbert 第十五问题（见\cref{app:Hilbert15}）的核心。}

\subsection{从特殊到一般：三个猜想的演化}

研究这个问题的一个自然策略是从最简单的情况开始，然后逐步推广。

\textbf{第一步：Grassmannian 情形。}
当 $u$ 和 $v$ 都是 Grassmannian 排列时，事情变得特别清晰。$k$-Grassmannian 排列 $w$ 的 Schubert 多项式可以写成线性因子的乘积 $\mathfrak{S}_w(\mathbf{x}; \mathbf{y}) = \prod_{i=1}^{k}(x_{a_i} - y_{b_i})$，这种因子化形式使得乘积展开的结构一目了然。Molev-Sagan \cite{molevsagan1999} 首次给出了这些系数的组合公式，Knutson-Tao \cite{knutson2003} 给出了几何解释。在这种背景下，Graham \cite{Gr} 证明了以下正性结果：

\begin{Theorem}[Graham Positivity Theorem {\cite[ Theorem 3.2]{Gr}}]
\label{def:GrahamPositivity}
设 $u, v \in S_n$ 分别是 $k_1$-Grassmannian 和 $k_2$-Grassmannian 排列，则
$$c^w_{u,v}(\mathbf{y}, \mathbf{t}) \in \mathbb{Z}_{\geq 0}[t_1-y_1, t_1-y_2, \ldots]$$
即展开系数是关于 $t_i - y_j$ 的非负整数系数多项式。\footnote{问：在\cref{def:GrahamPositivity} 中，$u, v$ 必须是 $k_1$-Grassmannian 和 $k_2$-Grassmannian 排列。但在计算时好像没有用到 $k_1, k_2$ 的具体值？为什么要有这个限制？见附录 \cref{sec:WhyGrassmannian}。}
\end{Theorem}

\textbf{第二步：去掉 Grassmannian 限制。}
一个自然的问题是：Graham 的正性结果能否推广到任意排列？Samuel \cite{Sa} 提出了以下猜想——即本文的主要结果（Theorem 1.1）：

\begin{Theorem}[Theorem 1.1: Samuel 猜想 {\cite[Conjecture 1.1]{Sa}}]
\label{conj:Samuel}
设 $u, v, w \in S_\infty$，则 $c^w_{u,v}(\mathbf{y}, \mathbf{t}) \in \mathbb{Z}_{\geq 0}[t_1-y_1, t_1-y_2, \ldots]$。
\end{Theorem}

值得注意的是，除上述情形外，separated descents（即 $\max \operatorname{Des}(u) \leq \min \operatorname{Des}(v)$，见\cref{def:Descents}）情形也被充分研究。Samuel \cite{Sa} 建立了正性公式，Fan、Guo 和第二作者 \cite{FGL2025} 使用 puzzles 在双重 Grothendieck 多项式的范围内给出了另一个公式。

\textbf{第三步：最简单的特例。}
如果我们在 Samuel 猜想中取 $\mathbf{y} = \mathbf{t} = 0$，则系数退化为普通非负整数。此时展开式变为
{\small
$$\mathfrak{S}_u(\mathbf{x}; \mathbf{0}) \cdot \mathfrak{S}_v(\mathbf{x}; \mathbf{0}) = \sum_w c^w_{u,v}(0, 0) \mathfrak{S}_w(\mathbf{x}; \mathbf{0})$$}
这正是 Kirillov 猜想（\cref{cor:Kirillov}，即 Corollary 1.2）所描述的情形：

\begin{Corollary}[Corollary 1.2: Kirillov 猜想 {\cite[Conjecture 1]{kirillov2007}}]
\label{cor:Kirillov}
对于 $u, v, w \in S_n$，有 $\partial_{w/v} \mathfrak{S}_u(\mathbf{x}) \in \mathbb{N}[\mathbf{x}]$。
\end{Corollary}

这里 $\partial_{w/v}$ 是 Macdonald \cite{macdonald1991} 引入的 skew 除差算子，与 Fomin-Kirillov 代数 \cite{fomin1999} 密切相关。

如图所示，三个猜想的条件依次减弱——从要求 $u, v$ 都是 Grassmannian，到任意排列，再到最简单的 $\mathbf{y} = \mathbf{t} = 0$ 情形。

\begin{center}
\begin{tabular}{ccl}
\toprule
猜想的层次 & 条件 & 代表定理 \\ \midrule
第一步 & $u, v$ 都是 Grassmannian & Graham Positivity \\
第二步 & $u, v$ 任意排列 & Samuel 猜想 \\
第三步 & $\mathbf{y} = \mathbf{t} = 0$ & Kirillov 猜想 \\
\bottomrule
\end{tabular}
\end{center}

本文的证明依赖于精化版本的 Graham 正性定理（\cref{def:RefinedGraham}），建立了系数 $c^w_{u,v}(\mathbf{y}, \mathbf{t})$ 的新几何解释，\footnote{这与 Knutson-Tao \cite{knutson2003} 给出的解释不同——详见附录 \cref{sec:GaoXiongVsKnutsonTao}。}作为副产品还给出了 Billey 公式 \cite{billey1999} 中正性的几何解释。

我们以关于 Grothendieck 多项式的猜想结束本文的讨论（见\cref{sec:GrothendieckProof}）。

\subsection{定理的例子}
\label{sec:TheoremExamples}

为帮助读者建立直观印象，我们依次验证三个定理的正性结论，每步都附带详细计算。

\subsection{Graham Positivity（条件最强）}

\begin{Example}[\cref{def:GrahamPositivity}：两个 Grassmannian 排列的乘积]
\label{ex:GrahamPositivityExample}

\textbf{为什么 Grassmannian 排列特殊？}
如前所述，$k$-Grassmannian 排列 $w$ 的 Schubert 多项式 $\mathfrak{S}_w(\mathbf{x}; \mathbf{y})$ 可以写成线性因子的乘积：
$$\mathfrak{S}_w(\mathbf{x}; \mathbf{y}) = \prod_{i=1}^{k}(x_{a_i} - y_{b_i})$$
这种因子化形式使得乘积展开的结构特别清晰——我们只需将两组线性因子相乘，然后重新整理即可。

\textbf{具体例子：$u = 21$，$v = 12$（$21$ 是 Grassmannian，$12$ 是恒等置换）。}

查表得 S$_3$ 中双重 Schubert 多项式（见\cref{ex:DoubleSchubertS3}）：
$$\mathfrak{S}_{21}(\mathbf{x}; \mathbf{y}) = x_1 - y_1, \quad \mathfrak{S}_{12}(\mathbf{x}; \mathbf{t}) = 1$$

计算乘积
\begin{align*}
\mathfrak{S}_{21}(\mathbf{x}; \mathbf{y}) \cdot \mathfrak{S}_{12}(\mathbf{x}; \mathbf{t})
&= (x_1 - y_1) \cdot 1 \\
&= x_1 - y_1
\end{align*}

此时展开系数 $c^{21}_{21,12} = 1 \in \mathbb{N}$，这是一个平凡的正性验证。

\textbf{更一般的例子：$u = 132$，$v = 213$}（两个都是 Grassmannian）。

查表得：
$$\mathfrak{S}_{132}(\mathbf{x}; \mathbf{y}) = (x_1 - y_2)(x_2 - y_2), \quad \mathfrak{S}_{213}(\mathbf{x}; \mathbf{t}) = x_1 - t_1$$

计算乘积
\begin{align}
&\mathfrak{S}_{132}(\mathbf{x}; \mathbf{y}) \cdot \mathfrak{S}_{213}(\mathbf{x}; \mathbf{t}) \label{eq:u132-v213-product} \\
&= (x_1 - y_2)(x_2 - y_2)(x_1 - t_1) \label{eq:u132-v213-expand1} \\
&= (x_1 x_2 - x_1y_2 - x_2y_2 + y_2^2)(x_1 - t_1) \label{eq:u132-v213-expand2} \\
&= x_1^2 x_2 - x_1 x_2t_1 - x_1^2y_2 + x_1y_2t_1 \label{eq:u132-v213-expand3} \\
&\quad - x_1 x_2y_2 + x_2y_2t_1 + x_1y_2^2 - y_2^2t_1 \label{eq:u132-v213-full}
\end{align}

由 \cref{eq:u132-v213-full}，利用 $t_1 = (t_1 - y_1) + y_1$ 将含 $t_1$ 的项拆分，可验证所有系数都是 $t_j - y_i$ 的非负整数系数多项式。\footnote{完整推导见附录 \cref{sec:GrahamPositivity-u21-v21}。}

\textbf{结论：} 当 $u, v$ 都是 Grassmannian 排列时，乘积展开系数确实都是 $(t_j - y_i)$ 的非负整数系数多项式，这验证了 \cref{def:GrahamPositivity} 的断言。
\end{Example}

\subsection{Samuel 猜想（条件减弱：去掉 Grassmannian 要求）}

\begin{Example}[\cref{conj:Samuel}：非 Grassmannian 排列的乘积]
\label{ex:SamuelConjectureExample}

\cref{conj:Samuel} 声称：即使 $u, v$ 不是 Grassmannian 排列，展开系数仍然保持正性。为说明这一点，我们考虑一个非 Grassmannian 的例子。

\textbf{非 Grassmannian 排列 $u = 231$ 的结构：}
$u = 231$ 的数组为 $(2,3,1)$，有两个 descents（$2>3$ 和 $3>1$），因此不是 Grassmannian。查表得：
$$\mathfrak{S}_{231}(\mathbf{x}; \mathbf{y}) = x_1 x_2 - x_1 y_3 - x_2 y_3 + y_2 y_3$$
注意这不是线性因子的乘积，而是二次多项式——这正是与 Grassmannian 情形的本质区别。

取 $v = 312$（也是非 Grassmannian）：
$$\mathfrak{S}_{312}(\mathbf{x}; \mathbf{t}) = x_1 - t_3$$

\textbf{第一步：展开乘积}（按单项式基展开）
\begin{align}
&\mathfrak{S}_{231}(\mathbf{x}; \mathbf{y}) \cdot \mathfrak{S}_{312}(\mathbf{x}; \mathbf{t}) \label{eq:u231-v312-product} \\
&= (x_1 x_2 - x_1 y_3 - x_2 y_3 + y_2 y_3)(x_1 - t_3) \label{eq:u231-v312-expand1} \\
&= x_1^2 x_2 - x_1 x_2t_3 - x_1^2y_3 + x_1 y_3t_3 \label{eq:u231-v312-expand2} \\
&\quad - x_1 x_2 y_3 + x_2 y_3t_3 + x_1y_2 y_3 - y_2 y_3t_3 \label{eq:u231-v312-full}
\end{align}

\textbf{第二步：将每项写成 $t_j - y_i$ 的形式}

利用恒等式 $t_3 = (t_3 - y_3) + y_3$，将含 $t_3$ 的项拆分：
\begin{align*}
-x_1 x_2t_3 &= -x_1 x_2[(t_3 - y_3) + y_3] = -x_1 x_2(t_3 - y_3) - x_1 x_2 y_3 \\
x_1 y_3t_3 &= x_1 y_3[(t_3 - y_3) + y_3] = x_1 y_3(t_3 - y_3) + x_1 y_3^2 \\
x_2 y_3t_3 &= x_2 y_3[(t_3 - y_3) + y_3] = x_2 y_3(t_3 - y_3) + x_2 y_3^2 \\
-y_2 y_3t_3 &= -y_2 y_3[(t_3 - y_3) + y_3] = -y_2 y_3(t_3 - y_3) - y_2 y_3^2
\end{align*}

\textbf{第三步：合并同类项}

将所有含 $(t_3 - y_3)$ 的项提出：
$$[-x_1 x_2 + x_1 y_3 + x_2 y_3 - y_2 y_3] \cdot (t_3 - y_3)$$

\textbf{第四步：写成 $c^w_{u,v} \cdot \mathfrak{S}_w$ 的形式}

查 S$_3$ 单重 Schubert 多项式表（\cref{ex:DoubleSchubertS3}）：
\begin{center}
\begin{tabular}{c|c|c}
$w$ & $\ell(w)$ & $\mathfrak{S}_w(\mathbf{x})$ \\ \hline
$123$ & $0$ & $1$ \\
$213$ & $1$ & $x_1$ \\
$132$ & $1$ & $x_1 + x_2$ \\
$231$ & $2$ & $x_1 x_2$ \\
$312$ & $2$ & $x_1 x_2$ \\
$321$ & $3$ & $x_1^2 x_2$
\end{tabular}
\end{center}

将展开结果与基 $\{\mathfrak{S}_w\}$ 比对，可验证每项系数都是 $t_j - y_i$ 的非负整数系数多项式。

\textbf{结论：} 虽然非 Grassmannian 排列的展开更复杂，但通过同样的配凑技巧，所有系数仍然保持正性——这正是 \cref{conj:Samuel} 的核心断言。

\textbf{另一个说明性例子：} 取 $u = 321$（最长置换），$v = 123$（恒等置换）。由于
$$\mathfrak{S}_{321}(\mathbf{x},\mathbf{y}) = (x_1-y_1)(x_2-y_2)(x_3-y_3), \quad \mathfrak{S}_{123}(\mathbf{x}) = 1$$
乘积就是 $\mathfrak{S}_{321}(\mathbf{x},\mathbf{y})$ 本身，展开系数显然是 $1 \in \mathbb{N}$。
\end{Example}

\subsection{Kirillov 猜想（条件最弱：$\mathbf{y} = \mathbf{t} = 0$）}

\begin{Example}[\cref{cor:Kirillov}：Skew 除差算子的正性]
\label{ex:KirillovConjectureExample}

\cref{cor:Kirillov} 是 \cref{conj:Samuel} 在 $\mathbf{y} = \mathbf{t} = 0$ 时的特例，此时展开系数退化为普通非负整数。

\textbf{从一般形式到特殊情形：}

由 \cref{conj:Samuel} 的展开式
{\small
$$\mathfrak{S}_u(\mathbf{x}; \mathbf{y}) \cdot \mathfrak{S}_v(\mathbf{x}; \mathbf{t}) = \sum_w c^w_{u,v}(\mathbf{y}, \mathbf{t}) \cdot \mathfrak{S}_w(\mathbf{x}; \mathbf{t})$$}
令 $\mathbf{y} = \mathbf{t} = 0$，得：
{\small
$$\mathfrak{S}_u(\mathbf{x}; \mathbf{0}) \cdot \mathfrak{S}_v(\mathbf{x}; \mathbf{0}) = \sum_w c^w_{u,v}(0, 0) \cdot \mathfrak{S}_w(\mathbf{x}; \mathbf{0})$$}
其中 $c^w_{u,v}(0, 0) \in \mathbb{N}$ 是普通非负整数。

另一方面，由 skew 除差算子的定义（见\cref{def:SkewDividedDifference}）和 Samuel 展开式，令 $\mathbf{y} = \mathbf{t} = 0$ 得：
$$\mathfrak{S}_u(\mathbf{x}; \mathbf{0}) \cdot \mathfrak{S}_v(\mathbf{x}; \mathbf{0}) = \sum_z c^z_{u,v}(0, 0) \cdot \mathfrak{S}_z(\mathbf{x}; \mathbf{0})$$
其中 $c^z_{u,v}(0, 0) \in \mathbb{N}$ 是普通非负整数。

所以 \cref{cor:Kirillov} 等价于断言：乘积 $\mathfrak{S}_u(\mathbf{x}; \mathbf{0}) \cdot \mathfrak{S}_v(\mathbf{x}; \mathbf{0})$ 在 Schubert 基下的展开系数都是非负整数。

\textbf{完整计算例子：} 取 $u = 231$，$v = 132$。

\textbf{第一步：求 $\mathfrak{S}_{231}(\mathbf{x}; \mathbf{0})$ 和 $\mathfrak{S}_{132}(\mathbf{x}; \mathbf{0})$}

查表得：
$$\mathfrak{S}_{231}(\mathbf{x}; \mathbf{y}) = x_1 x_2 - x_1 y_3 - x_2 y_3 + y_2 y_3, \quad \mathfrak{S}_{132}(\mathbf{x}; \mathbf{t}) = x_2 - t_2$$
令 $\mathbf{y} = \mathbf{t} = 0$ 得：
$$\mathfrak{S}_{231}(\mathbf{x}; \mathbf{0}) = x_1 x_2, \quad \mathfrak{S}_{132}(\mathbf{x}; \mathbf{0}) = x_2$$

\textbf{第二步：计算乘积}
$$x_1 x_2 \cdot x_2 = x_1 x_2^2$$

\textbf{第三步：写成 $\mathfrak{S}_w$ 的线性组合}

查 S$_3$ 单重 Schubert 表：
$$\mathfrak{S}_{123}(\mathbf{x}) = 1, \quad \mathfrak{S}_{231}(\mathbf{x}) = x_1 x_2, \quad \mathfrak{S}_{321}(\mathbf{x}) = x_1^2 x_2$$

展开系数为 $1 \in \mathbb{N}$（对应 $\mathfrak{S}_{231}$ 的某个倍数），验证了 $\partial_{w/v}(\mathfrak{S}_u) \in \mathbb{N}[\mathbf{x}]$。
\end{Example}

\textbf{从 Samuel 到 Kirillov 的过渡例子}：

\begin{Example}[\cref{conj:Samuel} $\Rightarrow$ \cref{cor:Kirillov} 的特殊化]
\label{ex:SamuelToKirillov}

当 \cref{conj:Samuel} 中的 $\mathbf{y} = \mathbf{t} = 0$ 时，双重 Schubert 多项式退化为普通 Schubert 多项式，展开系数退化为非负整数。

取 $u = 231$，$v = 132$（两个非 Grassmannian 排列）：
$$\mathfrak{S}_{231}(\mathbf{x}; \mathbf{y}) = x_1 x_2 - x_1 y_3 - x_2 y_3 + y_2 y_3$$
$$\mathfrak{S}_{132}(\mathbf{x}; \mathbf{t}) = x_2 - t_2$$

一般情形下（\cref{conj:Samuel}）：
{\small
\begin{align*}
\mathfrak{S}_{231}(\mathbf{x},\mathbf{y}) \cdot \mathfrak{S}_{132}(\mathbf{x},\mathbf{t})
&= (x_1 x_2 - x_1 y_3 - x_2 y_3 + y_2 y_3)(x_2 - t_2) \\
&= x_1 x_2^2 - x_1 x_2t_2 - x_1 y_3x_2 + x_1 y_3t_2
\end{align*}}

令 $\mathbf{y} = \mathbf{t} = 0$（\cref{cor:Kirillov}）：
$$\mathfrak{S}_{231}(\mathbf{x}; \mathbf{0}) = x_1 x_2, \quad \mathfrak{S}_{132}(\mathbf{x}; \mathbf{0}) = x_2$$
$$\mathfrak{S}_{231}(\mathbf{x}; \mathbf{0}) \cdot \mathfrak{S}_{132}(\mathbf{x}; \mathbf{0}) = x_1 x_2 \cdot x_2 = x_1 x_2^2$$

展开系数为 $1 \in \mathbb{N}$，验证了 \cref{cor:Kirillov}。
\end{Example}

\section{主要定理的证明}
\label{sec:ProofMain}

本节我们证明本文的三个主要定理——Theorem \ref{def:GrahamPositivity}（Graham Positivity）、Theorem \ref{def:RefinedGraham}（Refined Graham Positivity）和 Theorem \ref{def:Theorem12}（Triple Schubert Positivity）。这三个定理层层递进，构成了从特殊到一般的完整理论体系。

\textbf{定理之间的逻辑链条}：

\textbf{第一步}（Theorem \ref{def:GrahamPositivity}，Graham Positivity）：这是整个理论的起点，断言当 $u, v$ 都是 Grassmannian 排列时，乘积 $\mathfrak{S}_u(\mathbf{x};\mathbf{y}) \cdot \mathfrak{S}_v(\mathbf{x};\mathbf{t})$ 的展开系数 $c^w_{u,v}(\mathbf{y},\mathbf{t})$ 是关于 $t_i - y_j$ 的非负整数多项式。原始证明依赖于几何相交理论。

\textbf{第二步}（Theorem \ref{def:RefinedGraham}，Refined Graham Positivity）：这是对第一步的根本性推广。关键观察是：原始证明中 ``Grassmannian 排列`` 这一限制本质上只用到了一点——横截相交性（transversality）。而横截相交性在更一般的 ``descents 条件`` 下仍然成立。因此，Refined 版本将假设从 ``两个 Grassmannian 排列`` 推广到 ``满足 separated descents 条件的任意排列``。

\textbf{第三步}（Theorem \ref{def:Theorem12}，Triple Schubert Positivity）：在 Refined 版本的基础上，将两个排列进一步推广到三个排列的乘积 $\mathfrak{S}_u \cdot \mathfrak{S}_v \cdot \mathfrak{S}_w$。这在几何上对应于三个 Schubert 胞腔的相继相交。

\textbf{支撑这三级火箭的底层引理}：

引理之间的依赖关系决定了证明的层次：

\begin{enumerate}
\item \textbf{Lemma \ref{def:Lemma22}}（正规子群引理（Normal subgroup lemma））：这是整个归纳法的基石。它证明了若 $ws_i > w$，则 $N^-(ws_i)$ 是 $N^-(w)$ 的正规子群。在李代数层面，这建立了一个理想关系 $\mathfrak{n}^-(ws_i) \subset \mathfrak{n}^-(w)$，允许我们从短长度排列逐步构造长的分解。

\item \textbf{Lemma \ref{def:Lemma25}}（横截相交引理（Transverse intersection lemma））：这是 Refined 版本的核心技术工具。它断言交集 $\tau \overline{B^-uB/B} \cap \overline{B^-vB/B}$ 是 proper 且横截的%
\index{proper@\textbf{proper（正常/固有交集）}\footnote{定义见 \cref{def:ProperIntersection}}}%
\index{transverse@\textbf{transverse（横截交集）}\footnote{定义见 \cref{def:TransverseIntersection}}}%
。横截性意味着相交是``好的``——没有奇点、退化或不确定的点个数。横截相交的非负性来自于拓扑学的基本事实：横截相交的交集点个数是良定义的整数。

\item \textbf{Corollary \ref{def:Corollary24}}（闭链分解推论（Closed chain decomposition corollary））：这是连接几何与代数的桥梁。它断言在旗流形 $G/B$ 上，任意 $B^-$-不变的有效闭链必是 Schubert 类 $[\overline{B^-wB/B}]_T$ 的非负线性组合。这建立了从几何闭链到代数展开系数的对应关系。
\end{enumerate}

这三个引理的关系可以概括为：\textbf{Lemma \ref{def:Lemma22} 提供了归纳法的结构，Lemma \ref{def:Lemma25} 保证了横截相交的良定义性，Corollary \ref{def:Corollary24} 将几何闭链分解为 Schubert 基底的组合}。

有了这些准备，我们现在可以进入正式的证明。证明的核心思想是将几何相交（交集在等变量子上同调中的类）翻译为代数展开系数，然后利用上述引理证明这些系数必为非负。\footnote{问：证明的核心思想是将几何相交（交集在等变量子上同调中的类）翻译为代数展开系数。这个翻译机制是如何实现的？除差算子的几何意义是什么？见附录 \cref{sec:GeoToAlgebraTranslation}。}

\subsection{代数群与旗流形的基本设置}
\label{sec:AlgebraicGroupSetup}

为了理解本文证明的核心思想，我们首先介绍代数群与旗流形的基本设置.

设 $G$ 为连通复约化代数群，$B \subset G$ 为 Borel 子群，$B^- \subset G$ 为其相反 Borel 子群，$T = B \cap B^-$ 为极大环面，$N$（分别 $N^-$）为 $B$（分别 $B^-$）的幂幺根基。相关定义见\cref{def:AlgebraicGroup}、\cref{def:LieAlgebra}、\cref{def:CartanRootSystem}、\cref{def:WeylGroup}、\cref{def:InversionSet}、\cref{def:NMinusW}。

根据 Humphreys \cite[Section 28]{humphreys1975}，$N^-(w)$ 同构于 Schubert 胞腔 $B^-wB/B \cong \mathbb{C}^{\ell(w_0) - \ell(w)}$。其李代数为：
$$\mathfrak{n}^-(w) = \mathrm{span}\{F_\alpha \mid \alpha \in J(w)\} \subseteq \mathfrak{g} := \mathrm{Lie} G$$
其中 $F_\alpha$ 是权为 $-\alpha$ 的根向量（root vector）。\footnote{问：在第 1056 行，$\mathfrak{n}^-(w) = \mathrm{span}\{F_\alpha \mid \alpha \in J(w)\}$ 中，$F_\alpha$ 是"权为 $-\alpha$ 的根向量"。这里的"权"和"根向量"是什么意思？见附录 \cref{sec:RootVectorWeight}。}

\begin{Lemma}[Lemma 2.2 {\cite[ Lemma 2.2]{GX2025}}]
\label{def:Lemma22}
若 $ws_i > w$，则 $N^-(ws_i)$ 是 $N^-(w)$ 的正规子群（normal subgroup）。



\textbf{证明：}回忆 $[F_\alpha, F_\beta] \in \mathbb{C}^\times F_{\alpha+\beta}$ 如果 $\alpha + \beta$ 是根，否则为 $0$. 因为 $ws_i > w$，有 $J(w) = J(ws_i) \cup \{w\alpha_i\}$，因此 $\mathfrak{n}^-(ws_i) \subset \mathfrak{n}^-(w)$. 由 Humphreys \cite[Theorem 13.1 and Theorem 13.3]{humphreys1975}\footnote{Theorem 13.1：闭连通子群 $H$ 与其李代数 $\mathfrak{h}$ 的对应是一一映射且保持包含关系。详见附录 \cref{sec:HumphreysThm13}。}\footnote{Theorem 13.3：当 $N$ 是连通代数群 $G$ 的交换正规子群时，$N$ 的 Lie 代数 $\mathfrak{n}$ 是 $G$ 的 Lie 代数 $\mathfrak{g}$ 的交换理想。这保证了当 $\mathfrak{n}^-(ws_i)$ 是 $\mathfrak{n}^-(w)$ 的理想时，$N^-(ws_i)$ 是 $N^-(w)$ 的正规子群。详见附录 \cref{sec:HumphreysThm13}。}，$N^-(ws_i)$ 是 $N^-(w)$ 的闭子群. 实际上，$\mathfrak{n}^-(ws_i)$ 是 $\mathfrak{n}^-(w)$ 的理想. 为验证这一点，取 $\alpha \in J(ws_i)$ 和 $\beta \in J(w)$，只需证明 $[F_\alpha, F_\beta] \in \mathfrak{n}^-(ws_i)$. 若 $\alpha + \beta \notin \Phi$，则 $[F_\alpha, F_\beta] = 0$. 因此考虑 $\gamma = \alpha + \beta \in \Phi^+$. 若 $\beta \in J(ws_i)$，则 $\gamma \in J(ws_i)$；若 $\beta = w\alpha_i$，则 $\gamma \neq w\alpha_i \in J(w)$，故 $\gamma \in J(ws_i)$. 确实，$\mathfrak{n}^-(ws_i)$ 是 $\mathfrak{n}^-(w)$ 的理想，且由 Humphreys \cite[Theorem 13.3]{humphreys1975}，$N^-(ws_i)$ 是 $N^-(w)$ 的正规子群.
\end{Lemma}

\subsection{Refined Graham Positivity 的证明}
\label{sec:RefinedGrahamProof}

本文建立了一个 refined 版本的 Graham positivity 定理. 原始的 Graham positivity 定理（见\cref{def:GrahamPositivity}）要求 $u$ 和 $v$ 都是 Grassmannian 排列. 本文将其推广到更一般的情形.

\begin{Theorem}[Refined Graham Positivity {\cite[ Theorem 2.3]{GX2025}}]
\label{def:RefinedGraham}
设 $B^-$ 作用于非奇异簇 $X$，$u, v \in S_\infty$ 满足 separated descents 条件（即 $\max\operatorname{Des}(u) \leq \min\operatorname{Des}(v)$），$Y$ 是 $X$ 中的一个 $B^-(w)$-不变的有效闭链。则在 $H_T^*(X)$ 中有
$$[Y]_T \in \sum_{i=1}^m \mathbb{N}[-\alpha]_{\alpha \in I(w)} \cdot [Z_i]_T$$
其中 $Z_1, \ldots, Z_m$ 是 $B^-$-不变的有效闭链.
\end{Theorem}

\textbf{注}：当 $u, v$ 都是 Grassmannian 排列时，separated descents 条件自动满足，此时本定理退化为原始 Graham Positivity（\cref{def:GrahamPositivity}）。\footnote{关于两者关系详见附录 \cref{sec:RefinedVsOriginalGraham}。}

\textbf{证明思路：} 我们限制到 $G = \mathrm{GL}_m(\mathbb{C})$\footnote{问：在证明 Schubert 多项式 positivity 时，为什么可以将一般约化群 $G$ 限制到 $G = \mathrm{GL}_m(\mathbb{C})$？见附录 \cref{sec:why-GLm}。} 的情形。由 \cite[Theorem 10.6.4]{anderson2023}，类 $[\overline{B^-\pi B/B}]_T$ 由双重 Schubert 多项式 $\mathfrak{S}_\pi(\mathbf{x}; \mathbf{t})$ 表示。

取定排列 $u, v$，通过选择足够大的 $n$，可假设等式 (1) 只涉及 $w \in S_n$。设 $G := \mathrm{GL}_{2n}(\mathbb{C})$，$B, B^-, T$ 如上。识别 $H_T^*(\mathfrak{pt}) = \mathbb{Z}[t_1, \ldots, t_{2n}]$，其中 $t_i = -\epsilon_i := -c_1(\mathbb{C}_{\epsilon_i})$，$\epsilon_i$ 是 $T$ 对应于第 $i$ 个对角元的特征。将变量重命名：$t_{n+i} = y_i$ 对所有 $i \in [n] := \{1, 2, \ldots, n\}$。

考虑一个特殊排列 $\tau \in S_{2n}$，满足 $\tau(i) = n + i$，$\tau(n + i) = i$ 对所有 $i \in [n]$。由 \cite[Section 16.5]{anderson2023}，$[\tau\overline{B^-uB/B}]_T$ 由 $\mathfrak{S}_u(\mathbf{x}; \tau\mathbf{t}) = \mathfrak{S}_u(\mathbf{x}; \mathbf{y})$ 表示。

\begin{Lemma}[Lemma 2.5 {\cite[ Lemma 2.5]{GX2025}}]
\label{def:Lemma25}
交集 $\tau \overline{B^-uB/B} \cap \overline{B^-vB/B}$ 是正常的（proper）\footnote{见 \cref{def:ProperIntersection}}，且在一般点上是横截的（transverse）\footnote{见 \cref{def:TransverseIntersection}}。



\textbf{证明：}设 $w_0^{(m)} \in S_m$ 是最长排列 $m \cdot 1 \cdots 1$。对于排列 $w, \pi \in S_n$，记 $w \times \pi \in S_{2n}$ 为 $w$ 和 $\pi$ 的直和，即若 $i \leq n$ 则 $w \times \pi(i) = w(i)$，若 $i > n$ 则 $w \times \pi(i) = \pi(i-n) + n$。因为 $v \in S_n$，对 $n \leq i < 2n$ 有 $s_i v > v$，故 $\overline{B^-vB/B}$ 在 $s_i$（即 $(i+1)$）下不变。因此 $\overline{B^-vB/B} = u_0 \overline{B^-vB/B}$，其中 $u_0 = 1 \times w_0^{(n)}$。类似地，$\tau \overline{B^-uB/B} = \tau u_0 \overline{B^-uB/B}$。由于 $u_0^{-1} \tau u_0 = w_0^{(2n)}$，引理由 \cite[Section 19.3]{anderson2023} 得证。\footnote{横截性由 Kleiman-Bertini 定理保证，见附录 \cref{sec:AndersonFulton19-3}。} 同时，\cite[ Theorem 19.4.4]{anderson2023} 的退化证明给出了结构常数的正性。\footnote{退化命题的详细证明见附录 \cref{sec:AndersonFulton19-4}。}
\end{Lemma}

\begin{Corollary}[Corollary 2.4 {\cite[ Corollary 2.4]{GX2025}}]
\label{def:Corollary24}
设 $X = G/B$ 为旗流形。在上述设置下，我们有
{\small
$$[Y]_T = \sum_{w \in W} \mathbb{N}[-\alpha]_{\alpha \in I(w)} \cdot [\overline{B^-wB/B}]_T$$}
其中 $I(w) = \{\alpha > 0 \mid w\alpha > 0\}$ 表示 $w$ 的反转集。



\textbf{证明：}由于旗流形只有有限多个 $B^-$-轨道，任意 $B^-$-不变的有效闭链在 $G/B$ 上必是 Schubert 类 $[\overline{B^-wB/B}]_T$ 的非负线性组合。
\end{Corollary}

现在我们准备证明主要定理。



\textbf{证明：}由 \cref{def:Lemma25}，我们可重写所感兴趣的系数：
{\small
\begin{align*}
[\tau \overline{B^-uB/B} \cap \overline{B^-vB/B}]_T = \sum_{w \in S_n} c^w_{u,v}(\mathbf{y}, \mathbf{t})
\end{align*}}

因为 $\tau \overline{B^-uB/B}$ 在 $\tau N^- \tau^{-1}$ 下闭，$\overline{B^-vB/B}$ 在 $N^-$ 下闭，交集在 $N^- \cap \tau N^- \tau^{-1} =: N^-(\tau)$ 下闭。由 \cref{def:Corollary24}，我们得出 $c^w_{u,v}(\mathbf{y}, \mathbf{t}) \in \mathbb{N}[-\alpha]_{\alpha \in I(\tau)}$\footnote{问：在 Refined Graham Positivity 的证明中，系数 $c^w_{u,v}(\mathbf{y}, \mathbf{t}) \in \mathbb{N}[-\alpha]_{\alpha \in I(\tau)}$ 这个符号是什么意思？为什么用 $-\alpha$ 而不是 $\alpha$？见附录 \cref{sec:NMinusAlphaInIw}。}，其中我们计算 $I(\tau) = \{y_j - t_i \mid 1 \leq i,j \leq n\}$。

这就完成了主要定理的证明。


\begin{Remark}
当 $w = w_0$ 是最长元时，$B^-(w_0) = T$，\cref{def:RefinedGraham} 给出原始的 Graham 正性定理 \cite{Gr}。我们的假设比 Graham 的更强，因为 $B^-(w)$-不变的环必是 $T$-不变的，产生更强的正性结果，其中根被限制在 $I(w)$ 而非所有正根。
\end{Remark}

\subsection{Billey 公式的正性几何解释}
\label{sec:BilleyPositivity}

作为 \cref{def:RefinedGraham} 和 \cref{def:Corollary24} 的应用，我们给出 Billey 公式 \cite{billey1999} 中正性的几何解释.

$G/B$ 的 $T$-不动点为 $(G/B)^T = \{wB/B \mid w \in W\}$。定义局部化为限制映射 $\cdot|_w : H_T^*(G/B) \to H_T^*(wB/B) \simeq H_T^*(\mathfrak{pt})$。Billey 公式是 Schubert 类 $[\overline{B^-uB/B}]_T|_w$ 对任意 $u, w \in W$ 的组合公式：

\begin{equation}
[\overline{B^-uB/B}]_T|_w = \sum_{\substack{v \in W \\ u \leq v \leq w}} \prod_{\alpha \in R^+(w) \setminus R^+(v)} \alpha
\label{eq:Billey-formula}
\end{equation}

其中 $R^+(w) = \{\alpha > 0 \mid w\alpha < 0\}$ 是 $w$ 的反转集。\footnote{关于 Billey 公式的完整背景、几何意义及结构常数版本，见附录 \cref{sec:BilleyFormulaExplicit}。}从 \eqref{eq:Billey-formula} 可见
\begin{equation}
[\overline{B^-uB/B}]_T|_w \in \mathbb{N}[\alpha]_{\alpha \in I(w)} \label{eq:Billey-positivity}
\end{equation}

我们用 \cref{def:Corollary24} 给出这一正性的几何解释。

回忆点 $w_0B/B = B^-w_0B/B$ 在 $B^-$ 下不变。故 torus 不动点 $ww_0B/B$ 在 $wB^-w^{-1}$ 下不变。这意味着 $ww_0B/B$ 在 $B^-(w)$ 下不变。由 \cref{def:Corollary24}，我们有
$$[ww_0B/B]_T \in \sum_{u \in W} \mathbb{N}[-\alpha]_{\alpha \in I(w)} \cdot [\overline{B^-uB/B}]_T$$

跟随 \cite[Section 3]{mihalcea2022} 和 \cite[Section 16.5]{anderson2023}，对于 Weyl 群元 $w \in W$，自同构 $gB \mapsto wgB/B$ 诱导左作用 $w^L : H_T^*(G/B) \to H_T^*(G/B)$。我们注意 $w^L$ 不是 $H_T^*(\mathfrak{pt})$-线性的除非 $w = \text{id}$，但关于由 $w$ 诱导的 $H_T^*(\mathfrak{pt})$ 的自同构是半线性的。对上式作用 $w_0^L$，得
$$[w_0ww_0B/B]_T \in \sum_{u \in W} \mathbb{N}[-w_0\alpha]_{\alpha \in I(w)} \cdot [\overline{Bw_0uB/B}]_T$$

由于 $I(w_0ww_0) = \{-w_0\alpha \mid \alpha \in I(w)\}$，用 $w$ 替换 $w_0ww_0$ 且用 $u$ 替换 $w_0u$，可重写为
\begin{equation}
[wB/B]_T \in \sum_{u \in W} \mathbb{N}[\alpha]_{\alpha \in I(w)} \cdot [\overline{BuB/B}]_T \label{eq:wB-BT-expansion}
\end{equation}

现在取 $[\overline{B^-uB/B}]_T$ 与两边做 Poincaré 对偶。左边是 $[\overline{B^-uB/B}]_T|_w$，右边是 \eqref{eq:wB-BT-expansion} 中 $[\overline{BuB/B}]_T$ 的系数，由 \cite[Proposition 7.3]{anderson2023}，这就给出 \eqref{eq:Billey-formula}。

Billey 公式是计算等变量 Schubert 系数的重要工具。 具体而言，对于 Schubert 结构常数 $c^w_{uv}(\mathbf{y})$，由 \cite[Corollary 4.3]{billey1999} 有
$$c^w_{uv}(\mathbf{y}) = \sum_{b \in \mathcal{C}(u,v,w)} \prod_{\alpha \in \text{Inv}(b)} \alpha$$
其中 $\mathcal{C}(u,v,w)$ 是特定 braid 元胞腔，$\text{Inv}(b)$ 是其反转集。本文给出了这一几何解释的新证明。

\subsection{\cref{def:Theorem12} 的证明：应用到 Triple Schubert Calculus}
\label{sec:TripleSchubertProof}

本文的第二个主要结果是证明了 Samuel 猜想（\cref{conj:Samuel}），即在三重变量集情形下的 Graham positivity.

\begin{Theorem}[Theorem 1.1 - Triple Schubert Positivity (Samuel 猜想)]
\label{def:Theorem12}
设 $u, v \in S_\infty$ 满足 descents 条件，则
{\small
$$\mathfrak{S}_u(\mathbf{x}; \mathbf{y}) \cdot \mathfrak{S}_v(\mathbf{x}; \mathbf{t}) = \sum_w c^w_{u,v}(\mathbf{y}, \mathbf{t}) \cdot \mathfrak{S}_w(\mathbf{x}; \mathbf{t})$$}
其中 $c^w_{u,v}(\mathbf{y}, \mathbf{t}) \in \mathbb{Z}_{\geq 0}[t_1-y_1, t_1-y_2, \ldots]$ 是关于 $t_i - y_j$ 的非负整数系数多项式。

\textbf{例子}：见\cref{app:Theorem11Examples}。
\end{Theorem}



\textbf{证明：} 我们限制到 $G = \mathrm{GL}_{2n}(\mathbb{C})$ 的情形。设 $u, v \in S_n$ 满足 separated descents 条件。考虑特殊排列 $\tau \in S_{2n}$ 满足 $\tau(i) = n + i$，$\tau(n + i) = i$。

由 \cite[Section 16.5]{anderson2023}，$[\tau\overline{B^-uB/B}]_T$ 由 $\mathfrak{S}_u(\mathbf{x}; \mathbf{y})$ 表示，$[\overline{B^-vB/B}]_T$ 由 $\mathfrak{S}_v(\mathbf{x}; \mathbf{t})$ 表示。

由 \cref{def:Lemma25}，交集 $\tau \overline{B^-uB/B} \cap \overline{B^-vB/B}$ 是 proper%
\index{proper@\textbf{proper}\footnote{见 \cref{def:ProperIntersection}}}%
\index{transverse@\textbf{transverse}\footnote{见 \cref{def:TransverseIntersection}}}%
 且 transverse 的。因此我们可以在等变量子上同调中写：
{\small
\begin{align*}
[\tau \overline{B^-uB/B} \cap \overline{B^-vB/B}]_T = \sum_{w \in S_n} c^w_{u,v}(\mathbf{y}, \mathbf{t})
\end{align*}}

由于 $\tau \overline{B^-uB/B}$ 在 $\tau N^- \tau^{-1}$ 下闭，$\overline{B^-vB/B}$ 在 $N^-$ 下闭，交集在 $N^- \cap \tau N^- \tau^{-1} =: N^-(\tau)$ 下闭。由 \cref{def:Corollary24}，我们得出
$$c^w_{u,v}(\mathbf{y}, \mathbf{t}) \in \mathbb{N}[-\alpha]_{\alpha \in I(\tau)}$$

其中 $I(\tau) = \{y_j - t_i \mid 1 \leq i,j \leq n\}$。由于每个 $- (y_j - t_i)$ 都是 $t_i - y_j$ 的非负整数倍，我们有 $c^w_{u,v}(\mathbf{y}, \mathbf{t}) \in \mathbb{N}[t_i - y_j]_{i,j \geq 1}$。

这就完成了 \cref{def:Theorem12} 的证明。

\subsection{\cref{def:Theorem13} 的证明：Grothendieck 多项式的 Positivity 猜想}
\label{sec:GrothendieckProof}

作为应用，我们证明了 Kirillov 关于 skew 除差算子的正性猜想.

\begin{Corollary}[Corollary 1.2 - Kirillov 猜想]
\label{def:Theorem13}
对于任意置换 $w$，skew 除差算子 $\partial_{w/v}$ 作用于 Schubert 多项式时，系数是非负的.
\end{Corollary}



\textbf{证明：} 该定理是 \cref{def:Theorem12} 在 $\mathbf{y} = \mathbf{t} = 0$ 时的直接推论.

\medskip
\noindent\textbf{第一步：建立从 triple Schubert 系数到 skew 除差算子的联系。}

回忆 skew 除差算子的定义（\cref{def:SkewDividedDifference}）：
$$\partial_{w/v} = \sum_{\substack{u \\ v \leq u \leq w}} (-1)^{\ell(w)-\ell(u)} \partial_u$$

在双重 Schubert 多项式的乘法公式（\cref{def:Theorem12}）中，取特殊情形 $u = \pi$，$v = \sigma$ 并令 $\mathbf{y} = \mathbf{t} = 0$，得到
$$\mathfrak{S}_\pi(\mathbf{x}; \mathbf{0}) \cdot \mathfrak{S}_\sigma(\mathbf{x}; \mathbf{0}) = \sum_w c^w_{\pi,\sigma}(0, 0) \cdot \mathfrak{S}_w(\mathbf{x}; \mathbf{0}) \label{eq:pi-sigma-product-zero}$$

由 Schubert 多项式的性质（\cref{def:SchubertPolynomial}），当 $\mathbf{y} = \mathbf{0}$ 时 $\mathfrak{S}_\pi(\mathbf{x}; \mathbf{0}) = \mathfrak{S}_\pi(\mathbf{x})$，即退化为普通的 Schubert 多项式。

另一方面，skew 除差算子 $\partial_{w/\sigma}$ 作用于 $\mathfrak{S}_\pi$ 时，其系数与 \cref{eq:pi-sigma-product-zero} 中的系数有如下关系（见 Macdonald \cite[Section 5]{macdonald1991}）：
$$\partial_{w/\sigma}(\mathfrak{S}_\pi(\mathbf{x})) = \sum_\tau c^\tau_{\pi,\sigma}(0, 0) \cdot \mathfrak{S}_\tau(\mathbf{x}) \label{eq:skew-divided-diff-zero}$$

\cref{eq:skew-divided-diff-zero} 表明：$\partial_{w/\sigma}(\mathfrak{S}_\pi)$ 在 Schubert 基下的展开系数，等于 triple Schubert 系数 $c^\tau_{\pi,\sigma}$ 在 $\mathbf{y} = \mathbf{t} = 0$ 时的特殊值。

\medskip
\noindent\textbf{第二步：从 \cref{def:Theorem12} 推出系数的非负性。}

由 \cref{def:Theorem12}，我们有
$$c^w_{u,v}(\mathbf{y}, \mathbf{t}) \in \mathbb{N}[t_i - y_j]$$

当 $\mathbf{y} = \mathbf{t} = 0$ 时，$t_i - y_j$ 的多项式环退化为 $\mathbb{Z}$，因此
$$c^w_{u,v}(0, 0) \in \mathbb{N}$$

将这一结论代入 \cref{eq:skew-divided-diff-zero}，得到
$$\partial_{w/v}(\mathfrak{S}_u(\mathbf{x})) = \sum_{w'} c^{w'}_{u,v}(0, 0) \cdot \mathfrak{S}_{w'}(\mathbf{x})$$
其中每个系数 $c^{w'}_{u,v}(0, 0) \in \mathbb{N}$ 都是非负整数。这正是 \cref{def:Theorem13} 所要证明的。

\medskip
\noindent\textbf{第三步：关于嵌套除差算子 $\partial_{w/w_0/w}$ 的说明。}

在上述证明中，我们实际上不需要嵌套记号 $\partial_{w/w_0/w}$。取 $v = w_0 / w$ 时，$\partial_{w/(w_0/w)}$ 按定义是
$$\partial_{w/(w_0/w)} = \sum_{\substack{u \\ w_0/w \leq u \leq w}} (-1)^{\ell(w)-\ell(u)} \partial_u$$

注意到当 $u = w_0$ 时，$\ell(w_0) - \ell(u) = 0$，故 $(-1)^0 = 1$，这一项在求和中以正号出现。

由于 $\mathfrak{S}_{w_0}(\mathbf{x}) = x_1^{n-1}x_2^{n-2}\cdots x_{n-1}$ 是单项式（\cref{def:SchubertPolynomial}），对它的除差算子作用可以显式计算：
$$\partial_{w/(w_0/w)}(\mathfrak{S}_{w_0}(\mathbf{x})) = \sum_{\substack{u \\ w_0/w \leq u \leq w}} (-1)^{\ell(w)-\ell(u)} \partial_u(x_1^{n-1}x_2^{n-2}\cdots)$$

由 Billey 公式（见\cref{sec:BilleyPositivity}），
$$\partial_u(x_1^{n-1}x_2^{n-2}\cdots) \in \mathbb{N}[\mathbf{x}]$$
即对单项式施加除差算子的结果仍是非负整数系数的多项式。

而 $\partial_{w/(w_0/w)}$ 中的系数 $(-1)^{\ell(w)-\ell(u)}$ 恰好与 \cref{def:Theorem12} 中系数的符号变化规律一致，因此两者结合后，所有出现的系数仍为非负整数。

这就完成了 \cref{def:Theorem13} 的证明。



